% SPDX-License-Identifier: CC-BY-4.0
% Copyright (C) 2026 Fabian Januszewski, Abhijit (the Langlandscape team)
\section{Results of the six stages}\label{sec:results}

\paragraph{How this section is made.} Each box below is a statement copied \emph{verbatim} from the chapter source
\texttt{mathematics/chapters/s0N/*.tex} by \texttt{tools/extract\_statements.py}, including its status markers and
\verb|\provisional| notes. Where the chapter prints the status in a paragraph directly after the statement (as in stage~02
and in parts of stage~04), that paragraph is copied as well, below the statement and marked as such; a \stformal{} there
is the Lean status of the result (for Thm~6.1 and Thm~6.3 a schematic interface), not a claim about its mathematics. The only changes are mechanical: the \verb|\label| is moved into the box header, and every
cross-reference is replaced by its number in \mainpdf{} (built from the chapters of this submission tree, 170 pp; Lean tree 2323a252); so ``Thm~21.5'' inside a box
means Theorem~21.5 of \mainpdf. Hypotheses such as (H0)--(H2), (P), (Lev) and the cited inputs are defined there.
Each header gives the number and page in \mainpdf{} (the page is that of the label; in a few cases the statement
starts on the next page). Proofs are not reproduced. A box may carry several status markers;
each marker refers to the part of the statement it is attached to. Table~\ref{tab:allresults} at the end of the section
lists every theorem, proposition, corollary and conjecture of the chapters with the markers occurring in its statement.

\subsection{Stage 01: quaternionic foundations}
Chapter: \mainpdf, pp.~1--24. Lean: 17 results of the chapter are marked formalized (milestone closure S01).
% SPDX-License-Identifier: CC-BY-4.0
% Copyright (C) 2026 Fabian Januszewski, Abhijit (the Langlandscape team)
% GENERATED by tools/extract_statements.py from mathematics/chapters/ -- do not edit by hand.

\begin{quoted}{Theorem}{2.1}{31}{thm:s01-mass}{Eichler mass formula; \texttt{if:s01-mass}}{chapters/s01/2-mass.tex}
\begin{enumerate}[label=(\alph*)]
\item $\displaystyle\sum_{[x]\in\XB{\hat\cO^\times}}\frac1{|\Gamma_x|}=\frac{1}{24}\prod_{\ell\mid D}(\ell-1)$
  (full-unit convention), equivalently $\sum 1/w_x=\varphi(D)/12$.
\item (level change) For open $U'\le U\le\hat\cO^\times$: $\sum_{X_B(U')}1/|\Gamma_{x'}(U')|=[U:U']\sum_{X_B(U)}1/|\Gamma_x(U)|$.
  With $w$ in place of $|\Gamma|$ the same holds iff $-1\in U'$ or $-1\notin U$; if $-1\in U\setminus U'$ then
  $2\sum_{X_B(U')}1/w_{x'}=[U:U']\sum_{X_B(U)}1/w_x$ (\lean{Langlandscape.Quat.massW\_level\_change\_two}).
\item $[\hat\cO^\times:\Uzero{N}]=\psi(N)=N\prod_{\ell\mid N}(1+\ell^{-1})$ and $[\hat\cO^\times:\Uone{N}]=\psi(N)\varphi(N)=N^2\prod_{\ell\mid N}(1-\ell^{-2})$.
\item Hence $\sum_{\XB{\Uzero{N}}}1/w_x=\varphi(D)\psi(N)/12$, and for $p\ge5$, $p\nmid D$, $s\ge1$, $(M,pD)=1$:
\[|S_s|=|\XB{\Uone{Mp^s}}|=\frac{\varphi(D)}{24}\,\psi(Mp^s)\,\varphi(Mp^s).\]
\end{enumerate}
Status: (a) \stcited{} \lit{voight\_2021\_quaternion-algebras}{Thm.~25.1.1, p.~415 (with $w_J=\#\cO_L(J)^\times/\{\pm1\}$,
hence $1/12$)}; \lit{eichler\_1955\_zur-zahlentheorie-der-quaternionen-algebren}{\S4 ``Klassenzahl und Maß'', p.~136ff.};
in Lean it is the registered hypothesis \pathtt{Hyp.EichlerMass} (\Tpath{Hyp/EichlerMass.lean}, $1/24$, full units), proved
for $D=2,11$ in tests. (b) \stformal{} \Tpath{P2/Mass.lean} \pathtt{mass\_orbit\_count}; the $w$-variant \stformal{} \pathtt{Quat.massW\_level\_change},
\pathtt{Quat.massW\_level\_change\_two} (orbit counting). (c) first half \stformal{} \pathtt{gl2\_borel\_index}, \pathtt{mass\_level\_change}; second half
\stproved{} below (count $|\GL_2(\Z/N)|/(N\varphi(N))$). (d) first half \stpmc{} from (a)--(c) (proof below; cf.\ \lit{voight\_2021\_quaternion-algebras}{Thm.~25.3.18,
p.~423}), \stpmc{} in Lean as \pathtt{eichlerOrder\_mass} (under \pathtt{Hyp.EichlerMass}); second half \stpmc{} from (a)--(c)
and Lemma~1.7(b) (all $\Gamma_x=1$), and \stformal{} under \pathtt{Hyp.EichlerMass} as \pathtt{Quat.card\_XB\_Utower}.
\end{quoted}

\begin{quoted}{Theorem}{6.1}{48}{thm:s01-JLclassical}{classical JL, Brandt form; \pathtt{if:s01-targets}; seed Thm.~3.1}{chapters/s01/6-targets.tex}
Let $(M,D)=1$; for the $\Gamma_0$ statement let $k\ge2$ be even. Choosing $\Qp\hookrightarrow\C$ (or working with the import's global splitting over a number
field), there is an isomorphism of modules over the algebra generated by the $T_\ell$, $\ell\nmid DM$,
\[\Brandt{\Uzero{M}}{V_k}\otimes\C\ \cong\ S_k(\Gamma_0(DM))^{D\text{-new}}\ \oplus\ \begin{cases}\C\cdot\mathbf 1&k=2,\\0&k>2,\end{cases}\]
where $\mathbf 1$ spans the functions factoring through $\Nrd$ (one-dimensional since $\Q_{>0}\backslash\hat\Q^\times/\Zhat^\times=1$),
and $H^1_{\rm par}(\Gamma_0(DM),V_k\otimes\C)\cong S_k\oplus\overline{S_k}$. At $\Gamma_1$-level, for $(N,D)=1$ and $k\ge2$ of any parity: for $k\ge3$,
$\Brandt{\Uone{N}}{V_k}\otimes\C\cong S_k(\Gamma_1(N)\cap\Gamma_0(D))^{D\text{-new}}$, and for $k=2$ the same after dividing
by the functions factoring through $\Nrd$, compatibly with $T_\ell,S_\ell$ for $\ell\nmid ND$.
\stcited{} \lit{jacquet-langlands\_1970\_automorphic-forms-on-gl2}{Thm.~16.1 (LNM~114, p.~494; p.~250 in the retypeset copy in \pathtt{literature/}) (representation-theoretic; JL call
their proof a sketch)}; \lit{buzzard\_2004\_p-adic-families-of-automorphic-forms\_preprint}{Thm.~2, p.~13 ($\Gamma_1$ form,
attributed to JL, Shimizu, Arthur)}; \lit{hijikata-pizer-shemanske\_1980\_basis-problem-for-modular-forms-announcement}{\S0,
p.~280 (basis problem, all $N$)}; \lit{shimizu\_1965\_zeta-functions-of-quaternion-algebras}{(scan; not text-verified)};
\lit{pizer\_1980\_an-algorithm-for-computing-modular-forms-on-gamma0-n\_scan}{(scan; not text-verified; algorithmic side)};
survey \lit{voight\_2021\_quaternion-algebras}{Remark~41.5.13, p.~778 (historical account of the correspondence; the
statement itself is not formulated there)}.
\par\smallskip\noindent{\scriptsize\itshape [Status paragraph following the statement in the chapter:]}\par
\noindent\emph{Lean status.} \lean{Langlandscape.Hyp.JLClassical} \stformal{} as a \emph{schematic interface} (S01-L13): only the case $k=2$
with trivial coefficients, the Eisenstein line $\C$ with $T_\ell=\ell+1$, and an abstract Hecke module in place of
$S_k(\Gamma_0(DM))^{D\text{-new}}$. It is a statement about the true space only once that space is substituted for the abstract datum, which cannot be checked
in Lean here; no later result may be advertised as ``under classical JL'' without this caveat.
\end{quoted}

\begin{quoted}{Theorem}{6.3}{49}{thm:s01-lambdaJL}{$\Lambda$-adic Hecke JL}{chapters/s01/6-targets.tex}
Let $\T^{\ord}(DM)$ be Hida's ordinary $\Lambda$-adic Hecke algebra of tame level $\Gamma_1(M)\cap\Gamma_0(D)$ and
$\T^{\ord}_B(M)$ the one acting on $\varprojlim_se\,\Brandt{U_s}{\cOE}$ (or on its homological dual). There is a surjection
$\T^{\ord}(DM)\twoheadrightarrow\T^{\ord}_B(M)$ with $T_\ell\mapsto T_\ell$ ($\ell\nmid DMp$), $U_\ell\mapsto U_\ell$ ($\ell\mid Mp$),
classical $\langle d\rangle\mapsto$ seed $\langle d\rangle^{-1}$ for $d$ prime to $Mp$ (the guard $\gcd(d,Mp)=1$, stated in \lean{Langlandscape.Hyp.LambdaJL} as \texttt{Nat.Coprime d (M*p)}, S04-L17a; the diamonds are defined only for such $d$; the seed's adelic diamond, right translation by $\mathrm{diag}(1,d)$; Remark~4.5, Lemma~5.5), and $U_\ell\mapsto+\Pi_\ell$ ($\ell\mid D$), $\Pi_\ell$ acting by right translation; Hsieh writes $-$, which is not imported (stage 05, \pathtt{rem:s05-hsieh}),
\provisional{F6/D3: sign fixed to $+\Pi_\ell$, Hsieh's sign declared not imported; before: $U_\ell\mapsto\pm\Pi_\ell$ ($\ell\mid D$; sign convention-dependent); the Hsieh citation below read ``$\mathrm{JL}(U_q)=-U_q$ for $q\mid N^-$'' without the remark that it is not imported.}
whose kernel cuts out the
$D$-new quotient: $\T^{\ord,D\text{-new}}\cong\T^{\ord}_B$.
\stcited{} \lit{longo-vigni\_2009\_quaternion-algebras-heegner-points-and-the-arithmetic-of-hida-families\_arxiv}{Prop.~6.1,
p.~27 (iso of $\tilde\Lambda$-algebras $\T^{\ord}_\infty\simeq\mathbb B^{\ord}_\infty$, with $T(n,n)=\langle n\rangle$ normalised as in
their [25])}; \lit{hsieh\_2017\_hida-families-and-p-adic-triple-product-l-functions\_arxiv}{\S4.5, p.~34 (surjection
$\mathrm{JL}:\T(N,\mathbb I)\to\T^D(N,\mathbb I)$, with his sign $\mathrm{JL}(U_q)=-U_q$ for $q\mid N^-$, not imported)}.
\par\smallskip\noindent{\scriptsize\itshape [Status paragraph following the statement in the chapter:]}\par
\noindent\emph{Lean status.} \lean{Langlandscape.Hyp.LambdaJL} \stformal{} as a \emph{schematic interface} (S01-L13): two conjuncts, the surjection
$\varphi$ with the generator matching (including classical $\langle d\rangle\mapsto$ seed $\langle d\rangle^{-1}$ and $U_\ell\mapsto+\Pi_\ell$ for $\ell\mid D$; the diamond field of the Lean hypothesis is corrected to the inverse diamond on the Lean side, S04) and $\ker\varphi=$ an abstract ideal cutting out the $D$-new
quotient. The $\GL_2$ side (algebra, generators, ideal) is abstract data, so the hypothesis is $\Lambda$-adic JL only after instantiation by the true objects.
\end{quoted}


\subsection{Stage 02: tree quotients}
Chapter: \mainpdf, pp.~24--50. 72 \verb|\lean| markers at the S02 closure.
% SPDX-License-Identifier: CC-BY-4.0
% Copyright (C) 2026 Fabian Januszewski, Abhijit (the Langlandscape team)
% GENERATED by tools/extract_statements.py from mathematics/chapters/ -- do not edit by hand.

\begin{quoted}{Theorem}{7.7}{54}{thm:s02-SA}{strong approximation for $B^1$ at $S=\{\infty,q\}$}{chapters/s02/2-quotient.tex}
Since $q\nmid D$, $B_q\cong M_2(\Qq)$ is split and $B^1_q\cong\SL_2(\Qq)$ is non-compact. Hence $B^1$ (embedded diagonally) is
dense in the restricted product $\Bhat^{1,(q)}:=\prod'_{\ell\ne q}B^1_\ell$ (with respect to the $\cO^1_\ell$).
\smallskip

\noindent\emph{Consequence (the shape used; \texttt{if:s02-SA}).} For every $q$-admissible $U$,
\[\Bhat^1=B^1\big(B^1_q\times U^{1,(q)}\big).\tag{SA$_U$}\]
\par\smallskip\noindent{\scriptsize\itshape [Status paragraph following the statement in the chapter:]}\par
\noindent\emph{Status.} The theorem is \stcited{}:
\lit{voight\_2021\_quaternion-algebras}{Main Thm.~28.5.3, p.~485}, applied to the global ring $R=\Z[1/q]$ with eligible set
$S=\{\infty,q\}$. In Voight's notation $\Bhat=\Bhat^S=\prod'_{v\notin S}B_v$ is the $S$-finite adele ring
(\lit{voight\_2021\_quaternion-algebras}{27.6.4, p.~469}; there, when $S$ is clear from the context, the superscript $S$ is dropped and replaced by hats), so his
$\Bhat^1$ is our $\Bhat^{1,(q)}$. The hypothesis ``$B$ is $S$-indefinite'' (\lit{voight\_2021\_quaternion-algebras}{Def.~28.5.1,
p.~485}: $S$ contains a place unramified in $B$) holds because $q\in S$ and $q\nmid D$.
Equivalently \lit{vigneras\_1980\_arithmetique-des-algebres-de-quaternions}{Ch.~III, Thm.~4.3, p.~81} ($H^1_KH^1_S$ dense in
$H^1_{\A}$ if $H^1_S$ is not compact; page-checked) and
\lit{platonov-rapinchuk\_1994\_algebraic-groups-and-number-theory\_2up}{\S7.4, Thm.~7.12, p.~427} (page-checked in the M8 litnote).
Kneser \lit{kneser\_1969\_lectures-on-galois-cohomology-of-classical-groups}{\S5.5, Sublemma~b, p.~94 (Tata pagination; p.~72 of the scan)} only states the result;
it is listed for orientation and not relied upon. The consequence (SA$_U$) is \stproved{} below from the theorem. It is
\lit{voight\_2021\_quaternion-algebras}{28.4.8, p.~484} (density implies $\Bhat^1\subseteq B^1\hat\cO^1$) and
\lit{voight\_2021\_quaternion-algebras}{Cor.~28.5.12, p.~486}; the shape (SA$_U$), with the factor $B^1_q$, is the projection away from $\infty$ of the \emph{second}
equality of (28.5.13) there (which multiplies in $B^1_S=B^1_\infty\times B^1_q$). Both are stated only for the norm-one units $\hat\cO^1$ of an
$R$-order $\cO$, not for a general open subgroup; the same proof extends, since it uses only that $\hat\cO^1$ is open (proof below); this matters because $U^{(q)}$ need not be the unit group of an order (e.g.\ $\Uone{N}$-levels).
In Lean, (SA$_U$) for all admissible $U$ is registered as the hypothesis $\HypSA$, \lean{Langlandscape.Hyp.SA} (shape
\lean{Langlandscape.Tree.SAat}, quantified over \pathtt{Tree.IsQAdmissible}); this is a registered hypothesis, not a proof, and the
theorem stays \stcited. The derivation of (SA$_U$) from density below is not formalised. Per level, $\HypSA$ is replaced by
Prop.~7.9.
\end{quoted}

\begin{quoted}{Proposition}{7.9}{55}{prop:s02-BG}{SA for one level as a finite check}{chapters/s02/2-quotient.tex}
Let $U$ be $q$-admissible. The maps
\begin{gather*}
\phi_0:\GU\backslash V_0(\cT_q)\to\XB U,\ \GU gv_0\mapsto[g];\\
\phi_1:\GU\backslash V_1(\cT_q)\to\XB U,\ \GU g\alpha v_0\mapsto[g\alpha];\\
\phi_E:\GU\backslash E(\cT_q)\to\XB{U'},\ \GU ge_0\mapsto[g]
\end{gather*}
($g\in B^1_q$) are well defined and \emph{injective}, and $\phi=(\phi_0,\phi_1,\phi_E):\GU\backslash\cT_q\to\BG(U,q)$ is a locally
surjective morphism of graphs. The following are equivalent:
\begin{enumerate}[label=(\roman*)]
\item $\phi_0$ is surjective;
\item (SA$_U$) holds and $\Xone(U)\to\XB{U}$ is surjective;
\item $\BG(U,q)$ is connected.
\end{enumerate}
Moreover:
\begin{enumerate}[label=(\alph*)]
\item (i) implies that $\phi_0,\phi_1,\phi_E$ are all bijective, so $\phi$ is an isomorphism of graphs $\GU\backslash\cT_q\cong\BG(U,q)$.
  This implication does not use the connectedness of $\cT_q$.
\item In all cases, $\phi$ is an isomorphism of $\GU\backslash\cT_q$ onto a union of connected components of $\BG(U,q)$ containing
  $([1],0)$; using that $\cT_q$ is connected, onto the connected component of $([1],0)$.
\end{enumerate}
\par\smallskip\noindent{\scriptsize\itshape [Status paragraph following the statement in the chapter:]}\par
\noindent\emph{Status.} \stproved{} (new in this plan). No strong approximation is used. The connectedness of $\cT_q$
(Lemma~7.3(a), \stcited) enters only in (i)$\Rightarrow$(iii) and in the last clause of (b); the implications
(iii)$\Rightarrow$(i)$\Leftrightarrow$(ii), the injectivity, local surjectivity, and (a) are tree-free.
In particular Thm.~7.7 together with Hasse--Schilling (Lemma~7.6(b)) implies (ii), hence (i) and the
bijections of (a) (tree-free), and hence (iii) (using connectedness of $\cT_q$). As hypotheses, SA alone does not give (ii).
Conversely, for an explicit $(B,U,q)$ the condition (iii) is a finite computation, and it discharges (SA$_U$) and the surjectivity
of $\Xone(U)\to\XB U$ for that level. Classically, SA enters at this point through the definite analogue of Ihara's lemma:
\lit{diamond-taylor\_1994\_non-optimal-levels-of-mod-l-modular-representations}{Lemma~2 and Cor.~1, pp.~445--446} show that $T_\ell-2$, $S_\ell-1$ annihilate $A[l]$, $A$ the kernel of
$(f_1,f_2)\mapsto f_1+\varphi_pf_2$ on $S_k(U;\Q_l/\Z_l)^2$, for almost all $\ell\equiv1\bmod d$ (Lemma~2; hypotheses $p\nmid N_U l$), whence the Eisenstein statement (Cor.~1), and the paper says this lemma ``plays the role of Ihara's result'' (p.~445).
\emph{Lean status.} \stformal{} (all tree-free parts; in the coset model of Lemma~7.3, with $\phi_0,\phi_1$ combined
into one map \texttt{phiV} on all vertex orbits to $\XB U\sqcup\XB U$):
\lean{Langlandscape.Tree.phiV\_injective}, \lean{Langlandscape.Tree.phiE\_injective},
\lean{Langlandscape.Tree.phiV\_oO}, \lean{Langlandscape.Tree.phiV\_tO} (Step~3),
\lean{Langlandscape.Tree.phi\_locSurj\_type0}, \lean{Langlandscape.Tree.phi\_locSurj\_type1} (Step~4);
(i) is \lean{Langlandscape.Tree.surjType0} (every class of $\XB U$ is $[g]$ with $g\in\SL_2(\Qq)$), and
\lean{Langlandscape.Tree.surjType0\_iff} identifies it with the surjectivity of $\phi_0$;
\lean{Langlandscape.Tree.surjType0\_of\_BGConnected} ((iii)$\Rightarrow$(i)),
\lean{Langlandscape.Tree.bijective\_of\_BGConnected} ((iii)$\Rightarrow$ bijections),
\lean{Langlandscape.Tree.bijective\_of\_surjType0} ((a)), \lean{Langlandscape.Tree.surj\_of\_SA} ((ii)$\Rightarrow$(i)),
\lean{Langlandscape.Tree.SA\_of\_surj} ((i)$\Rightarrow$(ii)). None of these takes a \texttt{Hyp.*} argument.
(i)$\Rightarrow$(iii) is \lean{Langlandscape.Tree.BGConnected\_of\_surj}, formalized modulo \texttt{Hyp.BTtree}; the ``in particular''
is \lean{Langlandscape.Tree.BGConnected\_of\_SA\_HS}, formalized modulo $\HypSA$, $\HypHS$, \texttt{Hyp.BTtree} (and
\lean{Langlandscape.Tree.surjType0\_of\_SA\_HS} for (i), modulo $\HypSA$, $\HypHS$ only). The last clause of (b) (image $=$ the
component of $([1],0)$) also uses connectedness of $\cT_q$; it has no separate Lean declaration (not formalised; it would take
\texttt{Hyp.BTtree}).
\emph{Hypothesis on $U$ in Lean} (audit note): all of these except \pathtt{BGConnected\_of\_SA\_HS} and \pathtt{surjType0\_of\_SA\_HS}
assume only that $U$ is a $q$-level (\texttt{IsQLevel}: $U\le\hat\cO^\times$ of finite index, factoring at $q$), not that $U$ is
$q$-admissible; the condition $\Nrd U=\Zhat^\times$ is not needed for the proposition, in line with the proof above (Step~2 uses only
$\Nrd U\subseteq\Zhat^\times$).
\emph{Deviation from the plan tex (refinement).} The tex's ``Moreover, in all cases the maps are injective and $\phi$ is an isomorphism
onto the component of $(1,0)$'' is split here into the injectivity/local surjectivity in the statement, (a) ((i)$\Rightarrow$ bijections,
tree-free) and (b) (union of components, tree-free; single component, with the tree).
\end{quoted}

\begin{quoted}{Proposition}{7.10}{57}{prop:s02-quot}{quotient graph; seed Prop.~quot}{chapters/s02/2-quotient.tex}
Let $U$ be $q$-admissible and assume (iii) of Prop.~7.9. By Prop.~7.9 this holds given
Thm.~7.7, Hasse--Schilling \emph{and} the connectedness of $\cT_q$ (Lemma~7.3(a), \stcited; \textbf{[tree]}):
$\HypSA\wedge\HypHS\wedge\mathrm{Hyp.BTtree}\Rightarrow$ (iii) (Lean \pathtt{Tree.BGConnected\_of\_SA\_HS}, which takes \texttt{Hyp.BTtree});
the bijections of (a) need only $\HypSA\wedge\HypHS$. As hypotheses, SA alone does not suffice.
(iii) itself is the working hypothesis here.
\begin{enumerate}[label=(\alph*)]
\item There are bijections
  \begin{gather*}
    \GU\backslash\{\text{type-}0\text{ vertices}\}\ \xrightarrow{\ gv_0\mapsto[g]\ }\ \XB{U},\\
    \GU\backslash\{\text{type-}1\text{ vertices}\}\ \xrightarrow{\ g\alpha v_0\mapsto[g\alpha]\ }\ \XB{U},\\
    \GU\backslash\{\text{edges}\}\ \xrightarrow{\ ge_0\mapsto[g]\ }\ \XB{\Uzero{q}U}
  \end{gather*}
  ($g\in B^1_q$). Under these bijections the endpoint maps $o,t$ become $s_0,s_1$. In other words,
  $\GU\backslash\cT_q\cong\BG(U,q)$ as graphs.
\item For $g\in B^1_q$, $\Stab_{\GU}(gv_0)=\Gamma_g$ (level $U$), $\Stab_{\GU}(g\alpha v_0)=\Gamma_{g\alpha}$
  (level $U$) and $\Stab_{\GU}(ge_0)=\Gamma_g$ (level $\Uzero{q}U$). For another representative
  $x'=\gamma xu$ one has $\Gamma_{x'}=\gamma\Gamma_x\gamma^{-1}$. Each $\Gamma_x$ has norm $1$
  (Lemma~1.7(a)), so it lies in $\GU$ after conjugation.
\item Let $T$ be any Hecke operator supported away from $q$ ($T_\ell,S_\ell,U_\ell$ for $\ell\mid M$,
  $\langle d\rangle$, $U_p$ with the normalisation $\NUp$ of stage 01, and $\Pi_\ell$ for $\ell\mid D$),
  acting compatibly on levels $U$ and $\Uzero{q}U$. Then $s_0^*,s_1^*,s_{0*},s_{1*}$ commute with $T$ on
  Brandt modules $\Brandt{\cdot}{V}$ and on their homological versions (on $\Brandt{\cdot}{V}$, $s_{j*}$ means the trace
  $\mathrm{tr}_j$ of Def.~7.13).
\end{enumerate}
\par\smallskip\noindent{\scriptsize\itshape [Status paragraph following the statement in the chapter:]}\par
\noindent\emph{Status.} (a): \stproved{} given (iii); \stpmc{} via Thm.~7.7 and Hasse--Schilling (through
(ii)$\Rightarrow$(i)$\Rightarrow$(a), which does not use connectedness of the tree). (b): \stproved; \emph{refinement of the plan tex} (which has ``given (iii)''): the three stabiliser identities
hold for every $g\in B^1_q$ without (iii); (iii) is used only to know that every class of $\XB U$, $\XB{U'}$ has such a tree
representative. (c): \stproved{} (unconditional).
\end{quoted}

\begin{quoted}{Proposition}{7.19}{66}{prop:s02-harm}{harmonic cochains; seed Prop.~harm}{chapters/s02/3-cohomology.tex}
Under the hypotheses of Prop.~7.18(b), let $c$ be a function on oriented edges with values in $V$
such that $c(\gamma e)=\gamma_p\,c(e)$ for $\gamma\in\GU$, $c(\bar e)=-c(e)$, and $\sum_{o(e)=v}c(e)=0$ for all
$v$. The $\GU$-invariant $V$-valued harmonic cochains of this kind are
\[\ker\Big(\Brandt{\Uzero{q}U}{V}\xrightarrow{(\mathrm{tr}_0,\,\mathrm{tr}_1)}\Brandt{U}{V}^{2}\Big)\ \cong\ H_1(\GU,V),\]
with the trace maps $\mathrm{tr}_j$ (the plan's ``$s_{j*}$'') of Def.~7.13. If $-1\in U$, assume $-1$ acts trivially on $V$
(the standing hypothesis of Def.~7.13, needed for $\nu$); otherwise replace $V$ by $V^+$, as $V^-$ contributes $0$ to both sides
(Remark~7.15). This is the \emph{$q$-new} submodule of the Brandt
module of level $\Uzero{q}U$; the term is defined here as this kernel. The isomorphism is $\nu_{U'}^{-1}$ followed by
Prop.~7.18(b), $H_1(\GU,V)=\ker\partial$.
\par\smallskip\noindent{\scriptsize\itshape [Status paragraph following the statement in the chapter:]}\par
\noindent The description of the harmonic cochains as $\ker(\mathrm{tr}_0,\mathrm{tr}_1)$ is \stproved{} given (iii); the isomorphism with
$H_1(\GU,V)$ is \stpmc{} given (iii) (via Prop.~7.18(b); registered in Lean as hypothesis \pathtt{Hyp.SerreCoh}). (In
the sum, ``$o(e)=v$'' runs over all oriented edges with origin $v$, in both orientation classes.)
\emph{Lean status.} For trivial coefficients $V=R$ the $q$-new module $\ker(\mathrm{tr}_0,\mathrm{tr}_1)$ is
\lean{Langlandscape.Tree.qNew}, and $\nu_{U'}:\ker\partial\cong\ker(\mathrm{tr}_0,\mathrm{tr}_1)$ is \stformal{}
\lean{Langlandscape.Tree.qNewEquivKerPartial} (also \lean{Langlandscape.Tree.normMap\_restrict\_bijective}), assuming all $w_x$
at level $U$ are units in $R$ (remark below); no (iii) is needed for it. The identification with harmonic cochains on the tree is not
formalised.
\end{quoted}


\subsection{Stage 03: weighted graphs, the definite Ribet sequence and the Ihara lemma}
Chapter: \mainpdf, pp.~50--82.
% SPDX-License-Identifier: CC-BY-4.0
% Copyright (C) 2026 Fabian Januszewski, Abhijit (the Langlandscape team)
% GENERATED by tools/extract_statements.py from mathematics/chapters/ -- do not edit by hand.

\begin{quoted}{Proposition}{8.8}{79}{prop:s03-mazur}{Mazur check; seed Example mazur}{chapters/s03/1-graphs.tex}
For every prime $q$,
\[\Phicomb\big(\GDR_q(\hat\cO^\times)\big)\ \cong\ \Z\big/\mathrm{num}\tfrac{q-1}{12},\]
and this group is cyclic, generated via $\bar\lambda$ of Lemma~8.4. For $q\in\{2,3\}$ both sides are trivial.
\stpmc: Lemma~8.7, Prop.~8.3, Lemma~8.4. Lean (S03-L5, L6):
\lean{Langlandscape.Graph.mazurNum}, \lean{Langlandscape.Graph.mazur\_prod\_mul\_sum}, \lean{Langlandscape.Graph.mazur\_weights\_pairwise\_coprime}
\stformal{} (abstract arithmetic core: for weights in $\{1,2,3\}$ with the counts $n_2,n_3$ and $\sum w^{-1}=\frac{q-1}{12}$ one has
$\prod w\cdot\sum w^{-1}=\mathrm{num}\frac{q-1}{12}$, and the weights are pairwise coprime; $q\ge5$ only); the identification of the weights with the DR graph
(\lean{Langlandscape.Ribet.sum\_inv\_wx\_Ohat}, \lean{Langlandscape.Ribet.mazur\_check}, \stformal{} modulo the registered hypothesis \lean{Langlandscape.Hyp.EllipticCount} (and \pathtt{Hyp.EichlerMass})); the cases $q=2,3$ are \stformal{} as well, through the $q\le3$ clause of the same hypothesis \lean{Langlandscape.Hyp.EllipticCount} (which asserts $|\XB{\hat\cO^\times}|=1$ there; \lean{Langlandscape.Ribet.mazur\_check} covers every prime $q$), so the chapter's elementary derivation of that case is not formalized
\end{quoted}

\begin{quoted}{Proposition}{9.5}{83}{prop:s03-phieis}{$\Phicomb$ is Eisenstein}{chapters/s03/2-eisenstein.tex}
Let $q$ be any prime, $U=\Uzero M$, $B=B_q$, $\ell\nmid qM$, and let $T_\ell$ act on $\Z[\XBq U]$ by the Brandt matrix (homological
convention, $[x]\mapsto\sum_i[xg_i]$). Then $T_\ell$ is self-adjoint for $\langle,\rangle_w$, so it acts on $\Phicomb(\GDR_q(U))$, and
$T_\ell-(1+\ell)$ annihilates $\Phicomb(\GDR_q(U))$. \stproved{} (proof below, steps (1)--(4); valid for all primes $q$, including $q\in\{2,3\}$, and all
$M$ with $q\nmid M$ and $\ell\nmid qM$; it uses Lemma~9.2, Prop.~3.3(a),(b), Prop.~8.3, and the
finiteness of $\XBq U$, Prop.~1.6, which makes $\GDR_q(U)$ a finite graph). Lean: abstract core \stformal{}: \lean{Langlandscape.Graph.heckeQ} and \lean{Langlandscape.Graph.phiComb\_eisenstein} (steps (2)--(4), abstract: an integer matrix $B$ with row sums $\ell+1$,
$w_{e'}B(e,e')=w_eB(e',e)$ and $w_e\mid B(e,e')$ off the diagonal, acting on $\Phicomb$ of the two-vertex lattice $\Z[E]^0$ with weights $w_e>0$). Instantiation at $U=\hat\cO^\times$: \lean{Langlandscape.Ribet.phiComb\_eisenstein\_Ohat} \stformal{} modulo the registered hypothesis (step (1)) \lean{Langlandscape.Hyp.BrandtStabFree}.

Cross-references (not needed for the proof): [Ri90, Thm~3.12, p.~449; proof pp.~449--450, Remark~3.13] \stcited{}
\lit{ribet\_1990\_modular-representations-of-gal-qbar-q-arising-from-modular-forms}{Thm 3.12, p.~449}. Ribet's $\Phi=\coker(\iota)$ is our
$\Phicomb$, and his proof is the divisibility of the off-diagonal Brandt entries $B_\ell(i,j)$ by $e(i)=w_i$, which is step~(1) below, a purely quaternionic
fact. The geometric analogue for $\Phi_{N,p}$ is [Ed91, Th.~1, p.~163] \stcited{}
\lit{asterisque\_1991\_courbes-modulaires-et-courbes-de-shimura\_vol}{Edixhoven, Th\'eor\`eme~1, p.~163}, which assumes $p>3$.
\end{quoted}

\begin{quoted}{Theorem}{10.2}{89}{thm:s03-ribet}{definite Ribet sequence as a vertex collapse; seed Thm ribet}{chapters/s03/3-ribet.tex}
Let $q\nmid D$, $U$ as in \texttt{if:s02-quot}, $\G:=\GU\quot\cT_q$, and $|\G|$ its underlying graph, with the orientation (C1).
\begin{enumerate}[label=(\alph*)]
\item Over $\Z$ there is an exact sequence, equivariant for all Hecke operators $T=[UgU]$ with $g$ supported at a finite set $\Sigma\not\ni q$ of
  places at which $U$ is factorisable (Prop.~7.10(c); this holds for all standard levels):
  \[0\to H_1(|\G|,\Z)\to\Z[\XB{U'}]^0\xrightarrow{\ \alpha=(-s_{0*},\,s_{1*})\ }\big(\Z[\XB U]^0\big)^{2}\to0.\]
\item Let $p\ge5$ and let $V$ be a finitely generated $\Zp$-module with continuous $U_p$-action. Then, with the traces of (C3),
  \[0\to H_1(\GU,V)\to\Brandt{U'}V\xrightarrow{\alpha}\Brandt UV^{2}\to H_0(\GU,V)\to0.\]
\item $\rk_\Z H_1(|\G|)=|\XB{U'}|-2|\XB U|+1$. At a level where all $\Gamma_x$ are trivial, this is $(q-1)|\XB U|+1$, as in
  \pathtt{if:s02-indef}. For $V=\Zp$, $p\ge5$: $H_1(\GU,\Zp)=H_1(|\G|,\Z)\otimes\Zp$.
\end{enumerate}
Part (a) holds over every commutative ring $R$ in place of $\Z$ (with $H_1(|\G|,R)$).
\emph{Status.} Given (iii): (a), (c) first two sentences \stproved{} (Lemma~10.1 and Prop.~7.10); (b) and the last
sentence of (c) \stpmc{} (Prop.~7.18(b), which rests on the cited tree property and Shapiro's lemma; Lean \pathtt{Hyp.SerreCoh}). Without
assuming (iii): \stpmc{} via $\HypSA\wedge\HypHS\wedge\mathrm{Hyp.BTtree}\Rightarrow$(iii) (Prop.~7.10). The definite form of Ribet's
surjectivity is [Ri90, Thm~3.15, p.~451]
\lit{ribet\_1990\_modular-representations-of-gal-qbar-q-arising-from-modular-forms}{Thm 3.15, p.~451} \stcited. Ribet proves it with Eichler's
theorems. [Ri90, Rem.~3.18, p.~452]
\lit{ribet\_1990\_modular-representations-of-gal-qbar-q-arising-from-modular-forms}{Rem.~3.18, p.~452} records the graph-connectivity
proof (Livn\'e, Edixhoven), which is ours. Lean (S03-L9, any \texttt{CommRing R}, under \pathtt{Tree.BGConnected}, or under $\HypSA\wedge\HypHS\wedge\mathrm{Hyp.BTtree}$ in the \texttt{\_of\_SA\_HS} versions): \lean{Langlandscape.Ribet.ribet\_exact}, \lean{Langlandscape.Ribet.ribet\_rank}, \lean{Langlandscape.Ribet.ribet\_hecke\_equivariant}, \lean{Langlandscape.Ribet.boundary\_BG\_eq}, \lean{Langlandscape.Ribet.boundary\_BG\_apply}, \lean{Langlandscape.Ribet.ribet\_exact\_of\_SA\_HS}, \lean{Langlandscape.Ribet.ribet\_rank\_of\_SA\_HS} \stformal{} for (a) and the rank in (c);
coefficient layer for (b), definitions and compatibilities only (S03-L14): \lean{Langlandscape.Ribet.CoeffData}, \lean{Langlandscape.Ribet.CoeffData.ofLocal},
\lean{Langlandscape.Ribet.alphaV}, \lean{Langlandscape.Ribet.alphaVee}, \lean{Langlandscape.Ribet.s0TrV\_comp\_s0PullV}, \lean{Langlandscape.Ribet.alphaV\_triv},
\lean{Langlandscape.Ribet.alphaVee\_triv} \stformal{}; the exactness statement (b) with coefficients $V$ itself is not formalized (\leanplanned)
\end{quoted}

\begin{quoted}{Corollary}{10.5}{90}{cor:s03-ihara}{definite Ihara lemma, (i)--(ii); seed Cor.~ihara}{chapters/s03/3-ribet.tex}
(i) $\alpha$ in Thm~10.2(a) is surjective over $\Z$.
(ii) For every abelian group $A$ (in particular $A=\F_\ell$ for every prime $\ell$), the kernel of
$\alpha^\vee:(A^{\XB U})^2\to A^{\XB{U'}}$, $(f_0,f_1)\mapsto s_1^*f_1-s_0^*f_0$, is the line of constants $A\cdot(\mathbf1,\mathbf1)$.
(ii$'$) Let $p\ge5$ and let $W$ be a $\Zp$-module with continuous $U_p$-action. Then
\[\ker\big(\alpha^\vee:\Brandt UW^2\to\Brandt{U'}W\big)=\big\{(f_w,\,f_{[q]^{-1}w}):w\in W^{U^1_p}\big\},\qquad f_w(x):=[\tau(x)]^{-1}w,\]
and this is $H^0(\GU,W)=W^{U^1_p}$.
\stproved{} for (i), (ii) given (iii) (Lemma~10.1 and Prop.~7.10). (ii$'$): \stpmc{} (proof below, new: it checks
$f_w\in\Brandt UW$ by (C5) and $s_1^*f_{[q]^{-1}w}=s_0^*f_w$, and uses Lemma~7.17 (given (iii); the tree property is cited)
and the density of $\GU$ in $U^1_p$, i.e.\ Thm.~7.7; see (R4)). Here $U$ is assumed factorisable at $p$, $U=U_p\times U^{(p)}$ (true
for all standard levels), so that $\det:U_p\to\Zp^\times$ is onto and (C6) applies; ``continuous'' refers to a Hausdorff topology on $W$ for which
$W$ is a topological $\Zp$-module (e.g.\ $W$ finitely generated with the $p$-adic topology, or discrete). Lean (S03-L9): \lean{Langlandscape.Ribet.ihara\_surj}, \lean{Langlandscape.Ribet.ihara\_surj\_of\_SA\_HS} (i), \lean{Langlandscape.Ribet.ihara\_ker\_const}, \lean{Langlandscape.Ribet.ihara\_ker\_const\_of\_SA\_HS} (ii, any $A$) \stformal{}; (ii$'$) is not formalized.
\end{quoted}

\begin{quoted}{Theorem}{11.3}{95}{thm:s03-iharaH0}{Ihara (iii): $H_0$ and its Hecke action}{chapters/s03/4-ihara.tex}
Let $p\ge5$ and let $V$ be a finitely generated $\Zp$-module with continuous $U_p$-action. Use conventions (C1)--(C6).
\begin{enumerate}[label=(\alph*)]
\item The map
  \[\varepsilon:\Brandt UV^2\to V_{U^1_p},\qquad \varepsilon(f_0,f_1)=\sum_{[x]\in\XB U}w_x^{-1}\big([\tau(x)]\,f_0(x)+[q\,\tau(x)]\,f_1(x)\big)\]
  is well defined and surjective, and $\ker\varepsilon=\mathrm{im}\,\alpha$. Hence $\coker\alpha\cong H_0(\GU,V)=V_{U^1_p}$.
\item For $g$ supported at $\ell\nmid Dpq$, and $U$ factorisable at $\ell$, $U=U_\ell\times U^{(\ell)}$ (true for the standard levels,
  Def.~1.2), the induced action of $T_g$ on $\coker\alpha$ is
  \[T_g\ \longmapsto\ \deg(UgU)\cdot[\ell^{\,v_\ell(\Nrd g)}].\]
  In particular, $T_\ell$ acts as $(1+\ell)[\ell]$ and $S_\ell$ acts as $[\ell]^2$, the exponent being $+1$. If $U_p$ contains
  the scalars, then $[\ell]^2$ is the scalar $\ell$, which acts on $\Vk$ as $\ell^{k-2}$, in agreement with the classical
  normalisation $S_\ell=\ell^{k-2}$. The exponent $+1$ is for the cohomological operators $T_g$; for the transposed (homological) operators of the normative convention of stage 02, \pathtt{if:s02-indef}, it is $-1$, and there $S_\ell$ acts as $[\ell]^{-2}$ (Remark~11.4,
  exponent table, (R3)).
\item On the $\chi$-eigenpiece of $V_{U^1_p}\otimes\cOE$ for a character $\chi:\Zp^\times\to\cOE^\times$ (via $[\cdot]$), $T_\ell$ acts by
  $(1+\ell)\chi(\ell)$.
\end{enumerate}
\emph{Status.} \stpmc{} (proof below; the stage tex has \sttoprove{} with a sketch, which upgrades the seed's ``Hecke action \stsketch''). The
well-definedness of $\varepsilon$, $\varepsilon\circ\alpha=0$, surjectivity and the formula $\varepsilon\circ T_g=\deg(UgU)[\ell^m]\circ\varepsilon$ are \stproved{}
outright; $\ker\varepsilon\subseteq\mathrm{im}\,\alpha$ uses Cor.~10.5(ii$'$), hence (iii) and the density of $\GU$ in $U^1_p$ (full strong
approximation, (R4)); the identification with $H_0(\GU,V)$ is Prop.~7.18(b),(c). Lean (S03-L11): the definitions \lean{Langlandscape.Ribet.tau}, \lean{Langlandscape.Ribet.diamond}, \lean{Langlandscape.Ribet.epsilon} and the predicate
\lean{Langlandscape.Ribet.EpsWellDefined} \stformal{} (definitions only); clause (a) is the \emph{registered hypothesis} \lean{Langlandscape.Hyp.IharaH0} over
\pathtt{Ribet.alphaV} (an assumption, not a formalized theorem); clause (b) is
maths only here (the Hecke operators with coefficients, S04-L15a, now exist in Lean, but clause (b) itself is neither stated nor proved), \leanplanned
\end{quoted}

\begin{quoted}{Conjecture}{12.10}{107}{conj:s03-iharaGL2}{Ihara lemma for $\GL_2$, Betti form}{chapters/s03/5-amalgam.tex}
Let $\m$ be a maximal ideal of $\T^{(Mql)}$ (Lemma~12.7); call $\m$ \emph{Eisenstein} if there is $d\ge1$ such that $T_p-2\in\m$ for all but finitely
many primes $p\equiv1\bmod d$. This is the notion of [DT, p.~439--440]
\lit{diamond-taylor\_1994\_non-optimal-levels-of-mod-l-modular-representations}{p.~439--440}; DT's second condition $S_p-1\in\m$ is automatic for $p\equiv1\bmod dl$, since the
scalar $p$ acts as $p^{k-2}$ (Lemma~12.4(b)) and $l\in\m$. Compare Cor.~11.5 on the quaternionic side.
For non-Eisenstein $\m$, $(\coker\Psi)_{\m,\mathrm{tors}}=0$. Equivalently, $\Psi\otimes\F_l$ is injective after localisation at
$\m$, because the groups $H^1(\cdot,V)$ are torsion-free (so the Tor sequence of $\otimes\F_l$ has no extra term; $H^2(\Gamma,V/l)$ plays no role) and $\Psi$ is injective (Lemma~12.8). \stexp{} (consistent with C3: all torsion found is
Eisenstein, at $l=k+1$). The classical route uses Ihara's lemma [Ih75] (not in \pathtt{literature/}; \texttt{ihara\_1966} there is a different paper) or the congruence subgroup
property of $\SL_2(\Z[\frac1q])$ (\pathtt{literature/serre\_1970\_*}, scan, theorem not checked). The indefinite quaternionic analogue
is [MS19, Thm~1.1]
\lit{manning-shotton\_2019\_iharas-lemma-for-shimura-curves-over-totally-real-fields-via-patching\_arxiv}{Thm 1.1, p.~2} \stcited.
The ``equivalently'' needs $\ker\Psi=H^1(G,V)=0$, which is Lemma~12.8. Writing $0\to K\to A\xrightarrow\Psi B\to C\to0$ with
$A=H^1(\Gamma,V)^2$, $B=H^1(\Gamma_0(Mq),V)$ free over $\Z_{(l)}$ (Prop.~12.3), one has $\ker(\Psi\otimes\F_l)\supseteq K/l$ and
$\ker(\Psi\otimes\F_l)/(K/l)\cong C[l]$; so, as $K=0$, $\Psi\otimes\F_l$ is injective after localisation at $\m$ iff $C[l]_\m=0$ iff $C_{\m,\mathrm{tors}}=0$
(all maps are $\T^{(Mql)}$-linear by Lemma~12.7).
\end{quoted}


\subsection{Stage 04: the ordinary and \texorpdfstring{$\Lambda$}{Lambda}-adic passage}
Chapter: \mainpdf, pp.~82--101.
% SPDX-License-Identifier: CC-BY-4.0
% Copyright (C) 2026 Fabian Januszewski, Abhijit (the Langlandscape team)
% GENERATED by tools/extract_statements.py from mathematics/chapters/ -- do not edit by hand.

\begin{quoted}{Lemma}{13.2}{109}{lem:s04-ekills}{$e$ kills degree, vertex and component terms; seed Lemma ekills}{chapters/s04/ordinary.tex}
Let $s\ge1$, $V=\Zp$, $U$ of level prime to $p$.
\begin{enumerate}[label=(\alph*)]
\item $\deg:\Zp[X_s(U)]\to\Zp$ satisfies $\deg\circ U_p=p\cdot\deg$. Hence $e$ kills $\Zp$ (with $U_p=p$), $e\,\Zp[X_s]^0=e\,\Zp[X_s]$,
  and, for $q\nmid Dp$ and $U_q=\GL_2(\Zq)$, $e$ kills $H_0(\Gamma_{UU_s},\Zp)\cong\Zp$ (with the $U_p$-action induced from
  $\coker\partial$, \texttt{if:s02-coh}).
  \provisional{S04-M1: the $H_0$ clause uses that $\BG(UU_s,q)$ is connected (as in Thm.~lambdaribet, F7); the stage tex's proof
  says ``the connected graph'' but the statement does not list the hypothesis.} Mathematically the hypothesis is a theorem:
  $\Nrd(UU_s)=\hat\Z^\times$ and $B_q$ is split, so strong approximation for $B^1$ (Eichler, Kneser) gives connectivity; it is kept as
  a hypothesis because the Lean side carries it as \texttt{Hyp.SA}. (Also, $H_0(\Gamma,\Zp)\cong\Zp$ holds for any group; connectivity
  is needed for its identification with $\coker\partial$.)
\item ($B=B^{(q)}$, the definite algebra ramified exactly at $\{q,\infty\}$, written $B_q$ in stage~03, Def.~8.6; not
  the local algebra $B\otimes\Q_q$ of (N4); here $U=\Uone M$, so that $UU_s=\Uone{Mp^s}$ is a level of Def.~8.6, defined on
  $\hat B^{(q)\times}$ by the same local conditions; DR graph $\GDR_q(UU_s)$, \pathtt{if:s03-graph}.) Let $U_p$ act on the vertex module $\Zp[\{v_0,v_1\}]$ by $p$; then
  $\partial$ is $U_p$-equivariant, $e$ kills the vertex term, and $e\,\Xcomb_q(UU_s)\otimes\Zp=e\,\Zp[\XBq{UU_s}]$.
\item ($B=B^{(q)}$ as in (b).) $\Phicomb(\GDR_q(UU_s))\otimes\Zp\cong\Zp/|X_s|$ carries the action $\varphi\mapsto\varphi\circ U_p^*$ induced from
  $\Hom(X,\Z)$, which is multiplication by $p$; so $e(\Phicomb\otimes\Zp)=0$. This is a non-trivial statement: $v_p|X_s|\ge 2s-2$.
\item (Any $B$.) The Brandt pairing $\langle[x],[y]\rangle=\delta_{xy}$ on $\Zp[X_s]$ restricts to a perfect pairing
  \[e\,\Zp[X_s]\times e^*\Zp[X_s]\to\Zp,\]
  and $\nu$ identifies $e^*\Zp[X_s]$ with $e\,\Zp^{X_s}$ for the \emph{cohomological} $U_p$ ($f\mapsto\sum_jf(\cdot\,\varpi_j)$).
  Thus the perfect pairing is the evaluation pairing $e\,\Zp[X_s]\times e\,\Zp^{X_s}\to\Zp$. ($e$ and $e^*$ differ in general:
  Remark~13.3.)
\end{enumerate}
Here $U_p^*=T_{\varpi_p^{-1}}=\langle p\rangle_MU'_p$, where $U'_p=[UU_s\,\eta_p\,UU_s]$, $\eta_p=\iota_p^{-1}\mathrm{diag}(1,p)$, and
$\langle p\rangle_M$ is right translation by the central idele $t$ with $t_\ell=p$ for $\ell\ne p$ and $t_p=1$; for
$U^{(p)}=\Uone M$ this is the diamond $\langle d\rangle$ with $d\equiv p\bmod M$, $d\equiv1\bmod p^s$ (Lemma~4.3(d)).
\stproved{} below ((a)--(c) as in the seed, with the missing equivariance details; (d) with the identification of $e^*$; the
$H_0$ clause of (a) under the connectivity hypothesis). \emph{Lean} (for $\Zp$, $U=\Uone M$, $p\ge5$; the argument \texttt{hdec} is discharged by
\pathtt{padicOrdDecomp\_holds}): (a), degree part, \stformal: \lean{Langlandscape.Ordinary.deg\_comp\_UpHom},
\lean{Langlandscape.Ordinary.ekillsA\_tower} ($e\,\Zp[X_s]^0=e\,\Zp[X_s]$; abstract form \lean{Langlandscape.Ordinary.ordPart\_ker\_deg}); (d), perfect pairing,
\stformal: \lean{Langlandscape.Ordinary.perfect\_ordPart\_pairing\_tower\_padic}, \lean{Langlandscape.Ordinary.brandtPairing\_eq\_deltaPairing}. For (b), (c) only
the \emph{abstract} parts are formalized (\lean{Langlandscape.Ordinary.ordPart\_le\_ker\_of\_vertex}, \lean{Langlandscape.Ordinary.ordPart\_tensor\_phiComb\_eq\_bot}, with the
graph inputs and the identification $\Phicomb\cong\Z/|X_s|$ as hypotheses; the valuation bound $v_p|X_s|\ge2s-2$ modulo \pathtt{Hyp.EichlerMass} and a level-index input
\texttt{hlev}). Not formalized: the $H_0$ clause, the flatness identification $\Xcomb\otimes\Zp=\ker\partial_{\Zp}$, and the identification of $e^*$ with $e\,\Zp^{X_s}$
under $\nu$. So the lemma as a whole stays \stproved{} (tex).
\end{quoted}

\begin{quoted}{Theorem}{14.4}{114}{thm:s04-control}{$\Lambda$-freeness and control; seed Prop.~free (b); \texttt{if:s04-Mord}}{chapters/s04/control.tex}
Let $U$ be of level prime to $p$ and $r:=\rk_{\Zp}e\,\Zp[X_1(U)]$.
\begin{enumerate}[label=(\alph*)]
\item $\Mord(U)$ is a free $\Lam$-module of rank $r$, and for every $s\ge1$,
  $\Mord(U)\otimes_\Lam\cOE[\Delta_s]\xrightarrow{\sim}e\,\cOE[X_s(U)]$, compatibly with all Hecke operators.
\item (weights; $U=\Uone M$) For $k\ge2$, $\Mord(U)\otimes_\Lam\Lam/P^{\rm ad}_k$ is dual to $e\,\Brandt{U\Kp1}{\Vk\otimes\cOE}$ (Def.~16.1).
  On the component where $\mu_{p-1}$ acts by $\omega^{-(k-2)}$ it is dual to $e\,\Brandt{U\cdot\iota_p^{-1}\Uzero{p}_p}{\Vk\otimes\cOE}$, for every $k\ge2$.
  For $k\ge3$ the latter equals $e\,\Brandt{U}{\Vk\otimes\cOE}$ by $p$-stabilisation (\stexp{} in our normalisation).
  The first sentence is claimed integrally (over $\cOE$) only if $p\nmid\varphi(M)$; for $p\mid\varphi(M)$ it holds after inverting $p$,
  and its integral form is \stsketch{} (Remark~16.2(iv)). For a general level $U$ prime to $p$, (b) is \stexp{}
  (source: Hida 1988, Thm.~9.4, not checked here).
  \provisional{F4: (b) restated so that the character piece and the level match; before: ``For $k\ge2$ and $\chi$ the component on which $\mu_{p-1}$ acts by $\omega^{-(k-2)}$, $\Mord(U)^{(\chi)}\otimes_\Lam\Lam/P^{\rm ad}_k$ is dual to $e\,M(U\Kp1;\Vk\otimes\cOE)$; for $k\ge3$ one may replace $U\Kp1$ by $U\cdot\iota_p^{-1}\Uzero p_p$ (Hida's control at $\Gamma_0(p)$-level; expected)''; the left side was one $\mu_{p-1}$-eigenpiece and the right side the whole level-$U\Kp1$ module.}
\end{enumerate}
Status: (a) \stproved{} from Lemma~14.3 and Prop.~14.2 (choose compatible bases by Nakayama;
$\Lam=\varprojlim\cOE[\Delta_s]$); proof below, which uses only stage-01 results and elementary commutative algebra (raised from
\stpmc{} after the chapter audit S04-M9). Also \stcited{} \lit{hsieh\_2017\_hida-families-and-p-adic-triple-product-l-functions\_arxiv}{Thm.~4.2 and
its proof, p.~32: $\varprojlim_\alpha e\cO[X_\alpha]$ is $\Lam[\Delta_p]$-free, via [Hid88b, Cor.~10.1]} (Hsieh's ``Cor.~10.1'' is
\emph{Theorem}~10.1 in Hida's paper; there is no Cor.~10.1) and
\lit{longo-vigni\_2011\_control-theorems-for-quaternionic-hida-families\_arxiv}{Prop.~3.6 and its proof, p.~28: $\varprojlim_r
S_2(U_r,\cO_f)^{\rm ord}$ is $\Lam$-free by [Hid88b, Cor.~10.4]}.
(b) for $U=\Uone M$: \stcited{}, integrally if $p\nmid\varphi(M)$ (otherwise the integral form is \stsketch); general $U$: \stexp.
\lit{hsieh\_2017\_hida-families-and-p-adic-triple-product-l-functions\_arxiv}{Thm.~4.2(2), p.~32 (forms side; via
[Hid88b, Thm.~9.4])}; \lit{longo-vigni\_2011\_control-theorems-for-quaternionic-hida-families\_arxiv}{Thm.~3.7, p.~28 (eigen-components,
over a field)}; the dictionary to our $\NUp$, $\Vk$ and $P^{\rm ad}_k$ is Remark~16.2 (S04-M4), whose key
computation is Def.~16.1. The original source is
\lit{hida\_1988\_p-adic-hecke-algebras-for-gl2-over-totally-real-fields}{Thm.~8.6, p.~362 (weight independence of $\mathcal V^{\rm ord}$);
Thm.~9.4, p.~369 (control); Thm.~10.1, p.~370 (freeness over $\cO[[Z_a]]$ under $p$-divisibility); Cor.~10.4, p.~373 ($\Lam$-freeness,
$p>2$)}: the theorem headers and page numbers are located in the OCR text, but the statements are only partly legible, so their exact
hypotheses are taken from Hsieh and LV11.
\emph{Lean status of (a):} the finite-level part is \stformal{} for $\Zp$, $U=\Uone M$, $p\ge5$: control as surjectivity plus kernel
(\lean{Langlandscape.Ordinary.ord\_control\_finite}, \lean{Langlandscape.Ordinary.TowerHyp.mordTrans\_surjective},
\lean{Langlandscape.Ordinary.Mord.proj\_surjective}), freeness at level $s$ (\lean{Langlandscape.Ordinary.ordFree\_finite}). \emph{Not} formalized:
the constancy of the rank in $s$ (\pathtt{ordRank\_const}); the $\Lam$-adic half ($\Mord$ free of finite rank over $\Lam$ and the control
isomorphism for all $s$) is only stated, as the open hypothesis \lean{Langlandscape.Hyp.MordLambdaFree} (Lean-side status: open, S04-L14b). Thus
(a) as a whole stays \stproved{} (tex), with the Lean part \emph{formalized modulo \pathtt{Hyp.MordLambdaFree}} for the $\Lam$-adic half. The
Lean hypothesis \lean{Langlandscape.Hyp.MordLambdaFree} is thus the statement of a tex-proved theorem (not a merely cited one), taken as an argument because the data
$G_n=\Delta_{n+1}$, the reduction maps and the diamond representations are not yet in Lean; it has no named consumer in the stage-05/06 plans.
\emph{(b), Lean form (M9.3 sub-check, recorded):} \lean{Langlandscape.Hyp.HidaControl} is schematic, on abstract \emph{modules} $M_k$, $M_2$, $Q_k$ (the tex theorem is stated on
duals, $P^{\rm ad}_k$-quotients); it carries the guard $p\nmid\varphi(M)$ (M9.3) and $U_1(M)$-type levels only; its arguments \texttt{hk} and \texttt{h$\varphi$} are not used in the
statement (audit item O1 not applied).
\end{quoted}

\begin{quoted}{Theorem}{15.1}{116}{thm:s04-lambdaribet}{$\Lambda$-adic definite Ribet sequence; seed Thm.~lambdaribet; \pathtt{if:s04-lambdaribet}}{chapters/s04/lambdaribet.tex}
Let $q\nmid Dp$, $U$ of level prime to $pq$ with $U_q=\GL_2(\Zq)$ ($q$-admissible in the sense of stage 02). Assume $\BG(UU_s,q)$ is
connected for every $s\ge1$ (equivalently, Prop.~quot holds at all these levels). This follows from \texttt{Hyp.SA} $\wedge$
\texttt{Hyp.HS} (in Lean also \texttt{Hyp.BTtree}); it is not a finite check, and it does not imply $H_0=V_{U_p^1}$.
\provisional{F7: connectivity hypothesis made explicit; before: no such hypothesis was stated (only the implicit ``$q$-admissible'').}
Here $\Gamma_{UU_s}:=(\cO_{UU_s}[1/q])^1$ is the norm-one $\{q\}$-arithmetic group (as in Step~2 of the proof of Prop.~18.1); it does not
contain the central $q^{\Z}$. Put
\[\Xord_{Dq}(U):=\varprojlim_se\,H_1(\Gamma_{UU_s},\Zp)\otimes\cOE\]
(transition maps: corestriction, i.e.\ the restriction of $\pi_*$ on $\Zp[X^{(q)}_s]$-chains; over $\Zp$, $p\ge5$,
$H_1(\Gamma,\Zp)=H_1(\BG,\Z)\otimes\Zp$ by stage 02, D1). Then:
\begin{enumerate}[label=(\alph*)]
\item the sequence of $\Lam$-modules, equivariant for all Hecke operators supported away from $q$ (including $U_p$, $\langle d\rangle$,
  $\Pi_\ell$ for $\ell\mid D$),
  \[0\to\Xord_{Dq}(U)\to\Mord(\Uzero{q}U)\xrightarrow{\ \alpha=(-s_{0*},\,s_{1*})\ }\Mord(U)^{\oplus2}\to0\]
  is exact;
\item all three terms are $\Lam$-free, and $\rk_\Lam\Xord_{Dq}(U)=\rk_\Lam\Mord(\Uzero qU)-2\rk_\Lam\Mord(U)$;
\item at finite level $s\ge1$: $\rk_{\Zp}H_1(\Gamma_{UU_s},\Zp)=(q-1)|X_s(U)|+1$, and $\rk_{\Zp}e\,H_1=\rk e\Zp[X^{(q)}_{s}]-2\rk e\Zp[X_s]$,
  the ``$+1$'' being the degree line killed by $e$ (here $X^{(q)}_s$ is the level $\Uzero qUU_s$).
\end{enumerate}
Status: (a) \stproved{} below given \pathtt{if:s03-ribet}, Lemma~13.2(a) and $\NUp$; (b) \stproved{} via
Thm.~14.4(a); (c) \stproved{} (seed rank check; \pathtt{if:s02-indef}). \emph{Lean status:} \stformal{} for $E=\Q_p$, $U=\Uone M$, $p\ge5$, $q\ne p$:
the finite-level sequence of (a) and (c) at every $s\ge1$ (\lean{Langlandscape.Ordinary.ordRibet\_exact\_Utower\_uncond},
\lean{Langlandscape.Ordinary.rank\_ordRibet\_uncond}, \lean{Langlandscape.Ordinary.rank\_ordRibet\_full\_uncond}; connectivity of $\BG(UU_s,q)$ is an explicit
argument, \pathtt{Hyp.PadicOrdDecomp} is discharged). The limit statement of (a) is \emph{formalized modulo \pathtt{Hyp.InvLimExact}}
(and the connectivity argument): \lean{Langlandscape.Ordinary.lim\_exact\_ribet} proves injectivity, exactness in the middle and surjectivity of $\alpha$ on the
limits, but takes the surjectivity of the $A$-transitions $e\,H_1(\Gamma_{UU_{s+1}})\to e\,H_1(\Gamma_{UU_s})$ as the literal argument
\texttt{hA : $\forall n$, Function.Surjective (H.transA n)} (not a term of type \lean{Langlandscape.Hyp.InvLimExact}: that hypothesis is the implication
``$\forall s\ge1$ \texttt{BGConnected} $\Rightarrow\forall n$, \texttt{Surjective (H.transA n)}'', so \texttt{hA} is its conclusion, available once the connectivity argument is supplied). This $A$-surjectivity is \emph{not} formalized (undischarged; the snake-lemma route was dropped in S04-L13b), so
(a) stays \stproved{} (tex) and is not upgraded. Not formalized: (b) (the $\Lam$-rank formula).
\end{quoted}

\begin{quoted}{Theorem}{16.6}{122}{thm:s04-Gor}{Gorenstein and freeness; \texttt{if:s04-Mord}, second half}{chapters/s04/hsieh.tex}
Let $p>3$, $p\nmid MD$, $N=MD$, and let $F\in eS(N,\chi,\mathbb I)$ be an ordinary $\Lam$-adic newform of tame conductor $N$ with
$\mathrm{supp}\,N^-\subseteq\Sigma^0_F$ (discrete series at every $\ell\mid D$; Hsieh's hypothesis for the factorisation of
$\lambda_F$ through JL, [Hs, \S4.5, p.~34]), satisfying (CR, $\mathrm{supp}\,N^-$) (Hyp.~16.5). Let $\m_F\subset\T^D(N,\mathbb I)$
be the maximal ideal containing $\ker\lambda^D_F$ (above the maximal ideal $\m$ of Hyp.~16.5; all localisations below are
at $\m_F$, of $\T^D(N,\mathbb I)$-modules). Then:
\begin{enumerate}[label=(\roman*)]
\item Hsieh's $eS^D(N,\mathbb I)_{\m_F}$ is free of rank one over $\T^D(N,\mathbb I)_{\m_F}$ (proof of [Hs, Thm.~4.5]);
\item $eS^D(N,\chi,\mathbb I)_{\m_F}$ is free of rank one over $\T^D(N,\chi,\mathbb I)_{\m_F}$, and $\T^D(N,\chi,\mathbb I)_{\m_F}$ is
  Gorenstein (ibid.);
\item the eigenline $eS^D(N,\mathbb I)[\lambda^D_F]$ is free of rank one over $\mathbb I$ (its generator $F^D$: the primitive JL lift)
  ([Hs, Thm.~4.5]);
\item if moreover $p\nmid\varphi(M)$, then $\T^D(N,\mathbb I)_{\m_F}=\T^D(N,\chi,\mathbb I)_{\m_F}$ is Gorenstein, and
  $(\Mord(\Uone M)\otimes_\Lam\mathbb I)_{\m_F}$, which is the $\iota_2$-twisted $\mathbb I$-dual of $eS^D(N,\mathbb I)_{\m_F}$
  (Remark~16.3), is free of rank one over $\T^D(N,\mathbb I)_{\m_F}$.
\end{enumerate}
If $p\mid\varphi(M)$, Gorenstein-ness of the full $\T^D(N,\mathbb I)_{\m_F}$ and the homological freeness in (iv) are not cited and
not proved here (\stexp). For other $\m$ satisfying (CR$_D$) (old at $M$, or of smaller conductor) the conclusions are \stexp{}
(cf.\ [Wil95, Thm.~2.1], [Hel07, Cor.~8.11]; not checked here).
\par\smallskip\noindent{\scriptsize\itshape [Status paragraph following the statement in the chapter:]}\par
\noindent Registered in Lean as the hypothesis \lean{Langlandscape.Hyp.HsiehGor} (cited, not discharged; used by \lean{Langlandscape.Defect.Setting.gor\_new}, stage~6, which instantiates it over $\Lambda$).
\noindent\stcited{} (i)--(iii) \lit{hsieh\_2017\_hida-families-and-p-adic-triple-product-l-functions\_arxiv}{Thm.~4.5 and its proof, p.~34};
(iv) \stproved{} from (ii): if $p\nmid\varphi(M)$, the torsion of the diamond group, of order $(p-1)\varphi(M)$, is prime to $p$, so
(after enlarging $\cO$) its idempotents lie in $\T^D(N,\mathbb I)$ and $eS^D(N,\mathbb I)=\bigoplus_{\chi'}eS^D(N,\chi',\mathbb I)$; the
characters $\chi'$ have distinct reductions (orders prime to $p$), so exactly one idempotent, that of $\chi$, survives in the local ring
$\T^D(N,\mathbb I)_{\m_F}$, giving $\T^D(N,\mathbb I)_{\m_F}=\T^D(N,\chi,\mathbb I)_{\m_F}$; dualising a free rank-one module over a
Gorenstein ring gives a free rank-one module. The proof of Thm.~4.5 cites [Wil95, Thm.~2.1], [Hel07, Cor.~8.11, Rem.~8.12]
(\lit{helm\_2004\_maps-between-modular-jacobians-and-jacobians-of-shimura-curves\_arxiv}{Cor.~8.11 and Rem.~8.12, p.~31; journal
version p.~110, not checked on this node}) and [PW11, Prop.~6.4, 6.5]
(\lit{pollack-weston\_2006\_anticyclotomic-mu-invariants-of-modular-forms\_arxiv}{Prop.~6.4, p.~19; Prop.~6.5, p.~20}); the
freeness statement itself there is PW Thm.~6.2, p.~19, transported to the definite side by Prop.~6.5 (Prop.~6.4 is the self-pairing
statement, conditional on freeness). The inputs carry their own hypotheses (PW's CR; Helm's ``controllable'' $\m$; Wiles' Thm.~2.1);
that (CR, $\mathrm{supp}\,N^-$) suffices for this chain is Hsieh's claim and is not re-verified here. Gorenstein
follows from freeness plus the perfect pairing $B_N$. The freeness in (i), (ii) is in the \emph{proof}, not the statement, of
Thm.~4.5 (seed cites only ``Thm.~4.5''), and Hsieh states Gorenstein only for the $\chi$-piece $\T^D(N,\chi,\mathbb I)_{\m_F}$. Passing between $\Mord$ (homological) and $eS^D$ (forms) uses Hsieh (4.6) and
Thm.~4.2(1); dualising a free rank-one module over a Gorenstein ring is again free of rank one.
\end{quoted}

\begin{quoted}{Proposition}{17.1}{124}{prop:s04-etaprinciple}{$\eta$-principle; \pathtt{thm:s04-etaB}, \texttt{lem:s06-S}}{chapters/s04/etaprinciple.tex}
Let $A$ be a commutative ring and $S$ a commutative $A$-algebra that is finite free as an $A$-module. Let $M$ be an $S$-module with an
$S$-linear isomorphism $e:M\xrightarrow{\sim}S$, and let $B:M\times M\to A$ be an $A$-bilinear form such that
\begin{itemize}
\item $B$ is \emph{perfect}: the $A$-linear map $M\to\Hom_A(M,A)$, $m\mapsto B(m,\cdot)$, is bijective;
\item $B$ is \emph{$S$-equivariant}: $B(sm,n)=B(m,sn)$ for all $s\in S$, $m,n\in M$.
\end{itemize}
Let $\mu:S\to A$ be a homomorphism of $A$-algebras, and let $m_0\in M$ generate the eigenline
\[M[\mu]:=\{m\in M:\ sm=\mu(s)m\ \text{for all }s\in S\},\qquad\text{i.e.\ }M[\mu]=A\,m_0.\]
Then:
\begin{enumerate}[label=(\alph*)]
\item $S$ is Gorenstein with trace $\ell(s):=B(m_1,s\,m_1)$, where $m_1:=e^{-1}(1)$ (\lean{Langlandscape.CommAlg.EtaPrinciple.isGorensteinWith\_of\_pairing});
\item $\mathrm{congruenceIdeal}(\mu)=\big(B(m_0,m_0)\big)$, the principal ideal of $A$ generated by $B(m_0,m_0)$
  (\lean{Langlandscape.CommAlg.EtaPrinciple.congruenceIdeal\_eq\_span\_pairing}).
\end{enumerate}
\stproved{} below. \stformal{} (S06-L17, \pathtt{CommAlg/EtaPrinciple.lean}), in full generality of the statement and without registered hypotheses: \lean{Langlandscape.CommAlg.EtaPrinciple.isGorensteinWith\_of\_pairing} for (a) and \lean{Langlandscape.CommAlg.EtaPrinciple.congruenceIdeal\_eq\_span\_pairing} for (b); the trace in Lean is \texttt{traceOfPairing}, $\ell=B(e^{-1}(1),e^{-1}(\cdot))$.
\end{quoted}


\subsection{Stage 05: the comparison hypothesis and the simplicial shape}
Chapter: \mainpdf, pp.~101--134. The hypothesis $\HypX$ is never called known; in Lean it is split into
\texttt{Hyp.X.char}, \texttt{Hyp.X.iso}, \texttt{Hyp.X.gal} (\S\ref{sec:lean}). FD1 (the model for $\HypX$:
$\mu_N$-model with Picard functoriality) and FD2 (the sign $U_q\mapsto+\Pi_q$, \S\ref{sec:hsieh}) were decided by Fabian
on 15 November 2026; the chapter text below still carries the \verb|\provisional| notes recording the change.
% SPDX-License-Identifier: CC-BY-4.0
% Copyright (C) 2026 Fabian Januszewski, Abhijit (the Langlandscape team)
% GENERATED by tools/extract_statements.py from mathematics/chapters/ -- do not edit by hand.

\begin{quoted}{Hypothesis}{18.15}{139}{hyp:s05-X}{$\HypX(M;q)$; \texttt{if:s05-hypX}; seed Hyp.~8.1}{chapters/s05/hypX.tex}
\emph{Data.} A prime $p\ge5$, a prime $q\ne p$, an integer $M\ge1$ with $(M,pq)=1$; $B=B_q$; for $s\ge0$:
$\Gamma_s=\Gamma_1(Mp^s)\cap\Gamma_0(q)$, $X_s$ its compactified modular curve over $\Q$ in the $\mu_N$-model, $J_s=\Jac X_s$ with Picard
functoriality (decided, Fabian, 2026-11-15: FD1), $\Xq(J_s)$ the character group of the toric part of the special fibre at $q$ of its Néron model,
$S_s=\XBq{\Uone{Mp^s}}$, $\Z[S_s]^0$ the degree-zero part of the homological Brandt module, $\cH_s$ and
$\theta_s:\cH_s\to\End\Z[S_s]^0$ as in Def.~18.4 ($\langle d\rangle\mapsto\langle d\rangle^{-1}$, $U_p\mapsto N_{U_p}$ for $s\ge1$,
$T_p\mapsto T_p$ for $s=0$, $U_q\mapsto+\Pi_q$).

\emph{Statement.} There is a family of isomorphisms of abelian groups $\varphi_s:\Xq(J_s)\xrightarrow{\sim}\Z[S_s]^0$, $s\ge0$, such that:
\provisional{S05-M3: one family $(\varphi_s)_s$, as (X3) links $\varphi_s$ and $\varphi_{s'}$; the stage tex's Lean track writes
$\forall s\,\exists\varphi$ (Remark~18.16).}
\begin{enumerate}[label=(X\arabic*)]
\item \emph{(Hecke)} $\varphi_s\circ h=\theta_s(h)\circ\varphi_s$ for all $s\ge0$ and all $h\in\cH_s$;
\item \emph{(pairing)} $\langle\varphi_sx,\varphi_sy\rangle_w=u_s(x,y)$ for all $s\ge0$ and $x,y\in\Xq(J_s)$, where $u_s$ is Grothendieck's
  monodromy pairing, normalised to be positive definite (\pathtt{rem:s05-illusie}), and $\langle[e],[e']\rangle_w=w_e\delta_{ee'}$ with
  $w_e=|\Gamma_e/(\Gamma_e\cap\{\pm1\})|$ (\pathtt{if:s03-graph});
\item \emph{(naturality)} for $s'\ge s\ge0$, let $\beta_{s's}:\Xq(J_{s'})\to\Xq(J_s)$ be the map on character groups induced by the
  Picard pullback $J_s\to J_{s'}$ along the natural projection $X_{s'}\to X_s$ (contravariance of $\Xq$). Then
  $\varphi_s\circ\beta_{s's}=\pi_*\circ\varphi_{s'}$, with $\pi_*:\Z[S_{s'}]^0\to\Z[S_s]^0$ the pushforward along $\pi:S_{s'}\to S_s$. (Direction: the torus of the special fibre of a semistable Jacobian is
  $H^1(\Gamma,\Z)\otimes\mathbb G_m$ for the dual graph $\Gamma$ \lit{mazur\_1977\_modular-curves-and-the-eisenstein-ideal}{Appendix, \S3,
  p.~183, for $N=p$}, so $\Xq=H_1(\Gamma,\Z)$, and the map on characters induced by Picard pullback is dual to pullback on $H^1$, i.e.\
  the pushforward on $H_1$ of the dual graphs; this is why (X3) pairs $\beta_{s's}$ with $\pi_*$. \stsketch{} at $\Gamma_1$-level.)
\item \emph{(Galois)} for all $s\ge0$: on $T_pJ_s$, with $G_{\Q_q}$ and $\cH_s$ acting, the filtration $T_t\subset T^{I_q}\subset T_pJ_s$
  (Thm.~18.30(v), stated for any semistable abelian variety) has, as $\cH_s[G_{\Q_q}]$-modules, the graded pieces
  \[\begin{gathered}\text{bottom }T_t:\ \Hom(\Z[S_s]^0,\Zp(1))\ (\text{transpose action}),\\
  \text{middle }T^{I_q}/T_t:\ T_pJ_1(Mp^s)^{\oplus2}\ (\text{via the two degeneracies; for the operators away from $q$}),\\
  \text{top }T_pJ_s/T^{I_q}:\ \Z[S_s]^0\otimes\Zp\ (\text{action }h\mapsto\theta_s(h^\Ros)),\end{gathered}\]
  where $h^\Ros$ is the Rosati dual (for $T_\ell$: $\langle\ell\rangle^{-1}T_\ell$); for $h$ with $h^\Ros\notin\cH_s$ ($U_p$, and $U_\ell$ for
  $\ell\mid M$, whose Rosati duals involve $w_{Mp^s}$), $\theta_s(h^\Ros)$ \emph{means} the $\langle,\rangle_w$-adjoint $\theta_s(h)^*$ (the two readings
  agree where both are defined, Lemma~18.12(b)); the arithmetic $\Frob_q$ acts on the bottom piece as $q\,U_q$
  (equivalently $\sigma_q=U_q^{-1}$ on $\Xq(J_s)$, i.e.\ the Frobenius permutation of $S_s$ is right multiplication by $\Pi_q^{-1}$ under
  $\varphi_s$; this uses the $\mu_N$-model)
  (Frobenius clause in the $\mu_N$-model with Picard functoriality, as decided (Fabian, 2026-11-15, FD1); with the naive model it would read
  $q\langle q\rangle^{-1}U_q$, Remark~18.10);
  and the monodromy operator $N$ from the top to the bottom piece is $x\mapsto\langle x,\cdot\rangle_w$ (via (X2)). Transported to
  $H^1_{\rm par}(\Gamma_s,\Zp)(1)$ by the comparison of $T_pJ_s$ with Betti cohomology.
\end{enumerate}
$\HypXchar(M;q):=$ (X1)$\wedge$(X2)$\wedge$(X3); $\HypXgal(M;q):=$ (X4) (presupposing the family $\varphi$); $\HypX:=$ both. \sthyp.
At $s=0$ and $M\le2$, the $u$-adjoint of $\beta_{s'0}$ is compared with the transfer $\pi^!$ with the factor $2$ of Remark~18.3 (Lemma~18.12(c)).
\end{quoted}

\begin{quoted}{Proposition}{18.6}{128}{prop:s05-sign}{$\varepsilon=+1$; exponent $+1$}{chapters/s05/sign.tex}
\begin{enumerate}[label=(\alph*)]
\item ($\Gamma_0$-level, weight $2$.) Let $N=qM'$ with $q\nmid M'$, and identify $\Xq(J_0(N))\cong\Z[\Sigma(M')]^0$ (Ribet's
  Prop.~3.1; Thm.~18.30(ii)), and $\Sigma(M')\cong\XBq{\Uzero{M'}}$ by Deuring's correspondence as in Ribet \S3. Then $U_q=T_q$ acts
  on $\Xq(J_0(N))$ as the permutation $[x]\mapsto[x\Pi_q]$; so $\varepsilon=+1$. Since $\Pi_q^2$ acts trivially on
  $\XBq{\Uzero{M'}}$, this permutation equals $[x]\mapsto[x\Pi_q^{-1}]$ (the Frobenius permutation). \stpmc.
  At $\Gamma_1$-level, in the $\mu_N$-model of the Setting, the Frobenius permutation is $[x]\mapsto[x\Pi_q^{-1}]$ and differs from
  $[x]\mapsto[x\Pi_q]$ in general: \stpmc{} given (X1) and FD1 (it follows from (X1) and the four-row table, Remark~18.10, anchored at the naive/Albanese row, which is \stpmc{} conditional on FD1; Ribet--Stein's Lemma~3.17 agrees under a convention).
  (The model is fixed by the decision FD1 (Fabian, 2026-11-15): the Frobenius permutation at $\Gamma_1$-level ($\mu_N$-model, Picard functoriality) is $\Pi_q^{-1}$; the status \stpmc{} above is that of the derivation (given FD1), not of the model choice.)
\item (Consistency with stage 01.) For $D=q=11$, $M'=1$: $\Pi_{11}=\mathrm{id}$ on $\Z[\XB{\hat\cO^\times}]$, and $a_{11}(11a)=+1$
  (Example~5.6); so $\varepsilon=+1$. \stproved{} that $\Pi_{11}=\mathrm{id}$; $\varepsilon=+1$ then follows given the \stknown{} value $a_{11}=+1$ and the compatibility of
  $\theta$ with $U_{11}$ on the cuspidal line ((X1) at $\Gamma_0(11)$, i.e.\ (a), or JL at $11$; Example~5.6 is itself \stpmc).
\item (Exponent.) Let $N\ge1$ with $(N,q)=1$ and let $\cH^{\rm new}$ be the image of the weight-$2$ Hecke algebra of
  $\Gamma_1(N)\cap\Gamma_0(q)$ (generated by the $T_\ell$, $\ell\nmid Nq$, the $U_\ell$, $\ell\mid N$, the classical
  $\langle d\rangle$, $d\in(\Z/N)^\times$, and $U_q$) in the endomorphisms of the $q$-new subspace. Let
  $\theta:\cH^{\rm new}\to\End\Z[\XBq{\Uone N}]^0$ be a ring homomorphism with $\theta(\langle d\rangle)=\langle d\rangle^{-\delta}$ (adelic)
  and $\theta(U_q)=\varepsilon\Pi_q^{\,e}$, $\varepsilon,\delta,e\in\{\pm1\}$. If $\langle q\rangle^2\ne1$ on $\Z[\XBq{\Uone N}]^0\otimes\C$ (equivalently:
  some character $\chi$ of $(\Z/N)^\times$ with $\chi(q)\ne\pm1$ occurs in the adelic diamond action on it), then $e=\delta$. If
  $\langle q\rangle^2=1$ there (e.g.\ $N\le2$, or $q^2\equiv\pm1\bmod N$: only even $\chi$ occur, since $\langle-1\rangle=1$ on the class set,
  see the proof), the relation imposes nothing. \stpmc{}; checked numerically in T05-2, Remark~18.7(2).
  \provisional{S05-M1: compared with the stage tex, (c) constrains only $\theta(\langle d\rangle)$ and $\theta(U_q)$ (the other generators are
  arbitrary), and the condition on $\chi$ is read on the definite side; both are harmless generalisations (self-audit F8).}
\item (The diamond convention is $\delta=+1$, among the dictionaries considered here.) Let $\theta'=(\theta'_s)_{s\ge1}$ be either $(\theta_s)_s$ or $(\theta^\Ros_s)_s$
  (Prop.~18.18). If (X1)--(X3) of Hypothesis~18.15 hold for some family $(\varphi_s)$ with $\theta'$ in place of $\theta$,
  and $\rk_{\Zp}e\,\Zp[S_s]\ge2$ for some $s\ge1$, then $\theta'=\theta$; i.e.\ $\delta=+1$ among $\theta_s$, $\theta^\Ros_s$ and the dictionaries with $T_\ell\mapsto T_\ell$ \emph{as soon as some $\langle\ell\rangle^2\ne1$ on $\Z[S_s]^0$} (at $p=5$, $s=1$, $M=1$ only even $\chi$ with $\chi^2=1$ occur, so this qualifier is needed), hence $e=+1$ whenever (c) applies:
  $U_q\mapsto+\Pi_q$. The dictionaries with $\delta=-1$ and $T_\ell\mapsto T_\ell$ are excluded by the $T_\ell$-relation as soon as some
  $\langle\ell\rangle^2\ne1$ on $\Z[S_s]^0$ (Lemma~18.12(a)). This agrees with stage 01's diamond flag on forms (adelic
  $\langle d\rangle$ $=$ classical $\langle d\rangle^{-1}$, Remark~4.5). \stpmc{} (the input ``Picard $U_p$ commutes with the
  pullbacks'' is Lemma~18.9, \stpmc; batch-5 closure, previously \sttoprove). The hybrid dictionary $\langle d\rangle\mapsto\langle d\rangle$, $T_\ell\mapsto\langle\ell\rangle T_\ell=T_\ell^*$, $U_q\mapsto\Pi_q^{-1}$, $U_p\mapsto N_{U_p}$ ($\delta=-1$) satisfies the
  relations of Lemma~18.12(a), (b) and commutes with $\pi_*$, and is \emph{not} excluded by these arguments (nor by (X3)). Nor is it excluded by T05-2: at the levels computed there ($s=0$, no $U_p$) it
  coincides with $\theta^\Ros_0$, which matches the classical eigensystems as well as $\theta_0$ does (Remark~18.7(2); this is
  Prop.~18.18 at a single level). Excluding it \emph{modulo cited results}: under the hybrid dictionary the Brandt eigenvector attached to an eigenform $f$ carries the eigensystem of $\bar f$ away from $Mpq$, but $U_p\mapsto a_p(f)$ in place of $a_p(\bar f)=\overline{a_p(f)}$. The eigensystem matching away from $Mpq$ is Jacquet--Langlands
  (\lit{jacquet-langlands\_1970\_automorphic-forms-on-gl2}{Ch.~III, \S14, Thm.~14.4, p.~470 (printed pagination; PDF p.~479)}, with strong multiplicity one), which says nothing about $U_p$ at $p\mid$ level. The $U_p$-step needs \emph{local newform theory at $p$}
  (Casselman 1973, \emph{On some results of Atkin and Lehner}, Math.\ Ann.\ 201; not in \pathtt{literature/}, not checked): if $f$ is \emph{new at $p^s$}, the $U_1(p^s)$-fixed space of the local component $\pi_p$ is a line on which $U_p$ acts by $a_p(f)$, and $\bar f$ is again new at $p^s$ with $a_p(\bar f)=\overline{a_p(f)}$; the Brandt side has the same $\GL_2(\Q_p)$-module, so the $U_p$-eigenvalue on the line of $\bar f$ must be $\overline{a_p(f)}$. Hence the hybrid is excluded at every level $s\ge1$ carrying a form $f$ that is $q$-new, \emph{new at $p^s$} and has $a_p(f)\notin\R$ (this forces the $p$-part of the nebentypus to be primitive of conductor $p^s$: for trivial $p$-part $a_p$ is real). The hypothesis ``new at $p$'' cannot be dropped: for a $p$-stabilisation $f$ of a $p$-old form with $\chi(p)=1$, the $U_p$-eigenvalues on the $\bar f$-old space are $\{\bar\alpha,\bar\beta\}=\{p/\alpha,\bar\chi(p)\alpha\}\ni\alpha=a_p(f)$, so no contradiction arises, although $\alpha\notin\R$. Thm.~6.3 for the $\Lambda$-adic form is not re-checked here. \stpmc{} (cited: JL, strong multiplicity one, local newform theory; the existence of such an $f$ at some level is an extra condition we do not verify in general). All of (d) is a statement about the shape of
  $\HypX$, which itself stays \sthyp.
\end{enumerate}
Part (e) of the stage tex (higher weight, local Jacquet--Langlands for $\St\otimes\chi$) is written in the section on higher weight
(WP S05-M2).
\end{quoted}

\begin{quoted}{Proposition}{18.24}{150}{prop:s05-Xiso}{$\HypXiso$ holds non-canonically}{chapters/s05/subquot.tex}
Under the hypotheses of Def.~18.23, assume:
\begin{enumerate}[label=(\roman*)]
\item $H:=eH^1_{\rm par}(\Gamma_s,\Zp)^{\pm}_\m$ is free of rank one over $\T_\m$, and $\T_\m$ is $\Zp$-Gorenstein;
\item the hypotheses of Prop.~18.25 hold for $A=\Zp$, $\T_\m$, $\fa$ and $\fb:=\ker(\T_\m\to\T''_\m)$ ($\T''$ the $q$-old
  quotient): $\T_\m$ is finite free over $\Zp$ (automatic: finite and torsion-free over the PID $\Zp$), $\T'_\m$, $\T''_\m$ are torsion-free
  (automatic: images in endomorphism rings of lattices), $\fa\cap\fb=0$ and $(\T_\m)_K=(\T'_\m)_K\times(\T''_\m)_K$ (new and old
  eigensystems are disjoint, as the $U_q$-eigenvalues on the $q$-old part have absolute value $\sqrt q$ and those on the $q$-new part
  absolute value $1$);
\item $\T'_\m\cong\T_{B,\m}$ via $\theta_s$ (\pathtt{if:s01-targets}), and $e\,\Zp[S_s]_\m$ is free of rank one over $\T_{B,\m}$ with
  $\T_{B,\m}$ Gorenstein (16.8, [Hs, Thm.~4.5]).
\end{enumerate}
Then both sides of $\HypXiso(s)$ are free of rank one over $\T'_\m$, hence isomorphic (non-canonically). \stpmc.\\\emph{Lean:} the algebraic half (for $H$ free of rank one over a local $A$-Gorenstein $T$, with $T/\fa$ Gorenstein and the hypotheses of Prop.~18.25, $H[\fa]\cong T/\fa$) is \stformal{} (\pathtt{Comparison/XIso.lean}, sorry-free): \lean{Langlandscape.Comparison.xiso\_of\_gorenstein}; the verification of (i)--(iii) for the Hecke algebras of the modular curve and the $\Lam$-adic compatibility are not formalized.
Compatibility in $s$ (choosing compatible generators along the tower) additionally needs $\Lam$-adic freeness of $\varprojlim_sH$ over
the $\Lam$-adic $\T_\m$ (Hida, [Hi86a], not in \pathtt{literature/}) and of ${\Mord}_\m$ (16.8); \stsketch.
\emph{Freeness of rank one} is Hsieh's Thm.~4.5, which also assumes his hypothesis (CR) and the support condition on $N^-$ (to be carried explicitly in (i), (iii)); Helm's Cor.~8.11 and Rem.~8.12 are multiplicity-one (dimension) statements, entering via Hsieh's proof.\\
References: \lit{hsieh\_2017\_hida-families-and-p-adic-triple-product-l-functions\_arxiv}{Thm.~4.5, p.~34, whose proof cites
[Wil95, Thm.~2.1] and [Hel07, Cor.~8.11, Rem.~8.12]};
\lit{helm\_2004\_maps-between-modular-jacobians-and-jacobians-of-shimura-curves\_arxiv}{Cor.~8.11, Remark~8.12, p.~31}.
\end{quoted}

\begin{quoted}{Proposition}{18.25}{151}{prop:s05-subquot}{$q$-new sub- versus quotient lattice; seed Prop.~subquot; \pathtt{if:s05-subquot}}{chapters/s05/subquot.tex}
Let $A$ be a Noetherian local domain with fraction field $K$, $\T$ a local, finite free $A$-algebra, and $\fa,\fb\subseteq\T$ ideals with
$\T/\fa$, $\T/\fb$ $A$-torsion-free, $\fa\cap\fb=0$ and $\T_K=(\T/\fa)_K\times(\T/\fb)_K$ (equivalently, by \pathtt{NewOldData.split\_iff}:
there is $c\in A\setminus\{0\}$ with $c\in\fa+\fb$). Put $\T':=\T/\fa$, $\T'':=\T/\fb$ and $\cc_{\T'}:=(\fa+\fb)/\fa\subseteq\T'$. Let $H$ be free
of rank $r\ge0$ over $\T$, and $H[\fa]:=\{h\in H:\fa h=0\}$. Then
\[H[\fa]=\fb H\cong\fb^{\oplus r},\qquad H/\fa H\cong\T'^{\oplus r},\]
the natural map $H[\fa]\hookrightarrow H\twoheadrightarrow H/\fa H$ is injective, and its cokernel is $(\T'/\cc_{\T'})^{\oplus r}\cong(\T/(\fa+\fb))^{\oplus r}$.
In particular the cokernel is killed by $c$. \stproved{} (below), using only Lemma~D(a) ($\Ann_\T\fa=\fb$; \texttt{if:s06-lemD}; there
\stproved).
\emph{Lean:} \stformal{} (\pathtt{CommAlg/SubQuot.lean}, sorry-free, standard axioms only): \lean{Langlandscape.CommAlg.subquot\_kernel}, \lean{Langlandscape.CommAlg.subquot\_coker}.
\end{quoted}

\begin{quoted}{Proposition}{18.20}{144}{prop:s05-Xlambda}{$\Lam$-adic version; seed Prop.~Xlambda}{chapters/s05/hypX.tex}
Assume $\HypXchar(M;q)$. All limits are over $s\ge1$ (at $s=0$ there is no $U_p$ in $\cH_0$, hence no $e$; $s=0$ plays no role).
\begin{enumerate}[label=(\alph*)]
\item The maps $\beta_{s's}\otimes\Zp$ restrict to $e(\Xq(J_{s'})\otimes\Zp)\to e(\Xq(J_s)\otimes\Zp)$, and the $\varphi_s$ induce an isomorphism
  \[\Phi:\varprojlim_{s\ge1}e\big(\Xq(J_s)\otimes\Zp\big)\xrightarrow{\sim}\Mord(M),\]
  compatible with $T_\ell\mapsto T_\ell$ ($\ell\nmid Mpq$), $U_\ell\mapsto U_\ell$ ($\ell\mid M$), $U_p\mapsto U_p$, $U_q\mapsto\Pi_q$, the classical
  $\langle d\rangle\mapsto$ the adelic $\langle d\rangle^{-1}$ for $d\in(\Z/M)^\times\times\mu_{p-1}$ (embedded in $(\Z/Mp^s)^\times$ via $d\bmod M$ and the
  Teichmüller lift mod $p^s$; the finite group of \pathtt{def:s04-Mord}), and $\iota$-semilinear for $\Lam$: $\Phi(\lambda x)=\iota(\lambda)\Phi(x)$.
  In particular $\iota$ carries the seed's weight-$k$ point $P_k=\ker([u]\mapsto u^{k-2})$ to $P^{\rm ad}_k=\ker([u]\mapsto u^{2-k})$ of stage 04.
  \stproved{} given 13.4(a).
\item The left side of (a) is free over $\Lam$ of rank $r=\rk_{\Zp}e\,\Zp[S_1]$, and for $s\ge1$ its $\iota$-twisted specialisation
  $(\varprojlim e(\Xq(J_\bullet)\otimes\Zp))\otimes_{\Lam,\iota}\Zp[\Delta_s]$ is $e(\Xq(J_s)\otimes\Zp)$, Hecke-equivariantly
  ($\Delta_s$ the image of $1+p\Zp$ in $(\Z/p^s)^\times$). \stpmc{} (transport of Thm.~control(a) of stage 04, 16.8).
\item (Twisted pairings.) For $s\ge1$ let $E_s:=e(\Xq(J_s)\otimes\Zp)$ ($U_p$ is invertible on $E_s$), $w=w_{Mp^s}$, and
  \[[x,y]_s:=u_s(x,\,w\,U_p^{-s}y)\qquad(x,y\in E_s).\]
  \begin{enumerate}[label=(c\arabic*)]
  \item $[\,,\,]_s$ is perfect and $\cH_s$-equivariant: $[hx,y]_s=[x,hy]_s$ for $h\in\cH_s$. \stpmc{} (batch 6: from Lemma~18.19, (R2), (X1), (X2) and
    13.4(d)).
  \item \emph{(Expected; not proved.)} Through $\varphi_s$, $[\,,\,]_s$ corresponds to the twisted Brandt pairing of stage 04 (Hsieh's $B_{N,s}$,
    \texttt{def:s04-BN}), compatibly along the tower. \stopen{} (expected: a heuristic with no precise statement yet; the label \texttt{conjectural} is used only for want of a weaker label). This is not yet a precise statement: the normalisation of the comparison
    (which modules, which constant) is item (vi) of Remark~18.22, which lists the missing steps (i)--(vi). Nothing in this chapter
    uses (c2) as proved.
    \provisional{S05-M10: (c2) downgraded from \stsketch{} to \stconj/\stopen, after S05-R1 (batch-6 re-read, GAP: the norm twist was missing from
    the list). The list is completed in Remark~18.22; Lemma~18.21 locates the norm twist on the Brandt side.
    Downstream uses flagged: the $\Hom_\Lam(\Mord,\Lam)(1)$-form of (d2), the sentence after the proof, Remark~18.36, and the
    $\eta^B_F$ row of the table in \S18.4.1.}
  \end{enumerate}
  \provisional{batch 6 (E3): the seed and stage tex twist by $U_p^{-s}w_{p^s}$. Two changes, both corrections.
  (i) We use $w_{Mp^s}$, not $w_{p^s}$. Stage 04's $\tau^D_{Np^\alpha}$ has components at all primes of $M$ (\texttt{def:s04-BN}, translation (ii)).
  Also, for the partial $w_{p^s}$, Atkin--Li Prop.~1.3 gives $w_{p^s}U_\ell w_{p^s}^{-1}=\langle\ell\rangle_{p^s}^{-1}U_\ell$ for $\ell\mid M$. That is not
  $U_\ell^\Ros$ on the $\ell$-old part, where $U_\ell$ is not normal. So the (R1)-analogue that (c1) needs fails for $w_{p^s}$ when $M>1$ (\stsketch).
  (ii) We apply $U_p^{-s}$ first, then $w$. With the other order, equivariance for $h=U_p$ would need $U_p^\Ros$ to commute with $U_p$, which fails in general.
  The order agrees with $B_{N,\alpha}$, where $U_p^{-\alpha}$ is applied before translating by $\tau^D$.}
\item (Monodromy filtration of the ordinary tower.) Assume also $\HypXgal$. Let
  \[\Tord:=\varprojlim_{s\ge1}e_*\,T_pJ_s,\]
  the limit along the Albanese pushforwards $\pi_*:J_{s'}\to J_s$. Here $e_*$ is the ordinary projector of the Albanese $U_{p,*}=U_p^\Ros$, which commutes
  with $\pi_*$ (Lemma~18.9(c)). Each $h\in\cH_s$ acts through its Albanese version $h_*:=h^\Ros$, and $\Lam$ through the Albanese diamonds.
  This is Hida's convention
  \lit{hida\_1986\_galois-representations-into-gl2-zp-x-attached-to-ordinary-cusp-forms}{\S8, pp.~595--596: $J_r$ viewed as
  $\mathrm{Alb}(X_r)$, covariant Hecke action, limit along the covariant maps, $e$ attached to $T(p)$}.
  The filtrations $T_t\subset T^{I_q}\subset T_pJ_s$ (Thm.~18.30(v)) induce a $\Lam[G_{\Q_q}]$-stable filtration
  $\Tord_t\subset{\Tord}^{I}\subset\Tord$ with:
  \begin{enumerate}[label=(d\arabic*)]
  \item \emph{top:} $\Tord/{\Tord}^{I}\cong\varprojlim_se(\Xq(J_s)\otimes\Zp)$, the left side of (a); hence $\Tord/{\Tord}^I\cong\Mord(M)$ via $\Phi$.
    Under this, $h_*\mapsto\theta(h)$ for the operators $h$ listed in (a), and the $\Lam$-structure is $\iota$-semilinear as in (a). $\Frob_q$ acts as $\Pi_q^{-1}$
    (the Frobenius clauses of (d1), (d2) come from the Frobenius clause of (X4); FD1 is decided (Fabian, 2026-11-15): the model of $X_s$ is the $\mu_N$-model with Picard functoriality.).
  \item \emph{bottom:} $\Tord_t\cong\Mord(M)\otimes\Zp(1)$, with the same Hecke and $\Lam$-structure. $\Frob_q$ acts as $q\,\Pi_q^{-1}$.
  \item \emph{monodromy:} under (d1), (d2), the monodromy operator $\Tord/{\Tord}^I\to\Tord_t(-1)$ is $\pm$ the identity of $\Mord(M)$, with $+$ for the
    identifications fixed in the proof.
  \item \emph{middle:} ${\Tord}^I/\Tord_t\cong\Tord(M)^{\oplus2}$, Hida's ordinary Tate module at level $M$, via the two degeneracy maps.
    It is $\Lam$-free \lit{hida\_1986\_galois-representations-into-gl2-zp-x-attached-to-ordinary-cusp-forms}{Thm.~3.1(ii), p.~562
    ($V^{\rm ord}$ is $\Lam$-free of finite rank), with $J^{\rm ord}_\infty\cong V^{\rm ord}$, p.~596}. Hence $\Tord$ is $\Lam$-free.
  \end{enumerate}
  Status: (d1)--(d3) \stpmc{} (batch 6). The Hecke and transition structure of (d1), (d2) uses only $\HypXchar$; the Frobenius and $N$ use the
  corresponding clauses of (X4). (d4) \stsketch{} (missing step at the end of the proof).
  \provisional{batch 6 (E3): statement of (d) changed. Batch 5 found that the projector and the tower were not fixed, and that the pieces had been computed
  with the wrong (Picard) projector. Now:
  \begin{itemize}
  \item $\Tord$ is Hida's Albanese tower with $e_*$.
  \item The bottom piece is $\Mord(M)(1)$, identified canonically through the perfect $(e,e^*)$ Brandt pairing. It replaces the former
    $\Hom_\Lam(\Mord,\Lam)(1)$; identifying the two needs the $\Lam$-adic twisted pairing of (c2), which is \stopen{} (\stconj; Remark~18.22),
    so the $\Hom_\Lam$-form is not claimed.
  \item The top piece is $\Mord(M)$ with $h_*\mapsto\theta(h)$. This is the former ${\Mord}^\Ros$ read through $h_*=h^\Ros$: the same Rosati twist as in (X4).
  \item Hida's $\Lam$-freeness is now needed only for the middle piece, and its page is checked.
  \end{itemize}
  The status is raised from \stsketch{} to \stpmc{} for (d1)--(d3).}
\end{enumerate}
\par\smallskip\noindent{\scriptsize\itshape [Status paragraph following the statement in the chapter:]}\par
\noindent\emph{Lean status} (part (a), finite level only): \stformal{} \emph{conditional on} $\HypXchar$ (carried as the argument \texttt{hX : Hyp.X.char}), for an abstract \texttt{CharGroupData} $G$ and not for $\Xq(J_s)$ itself (\lean{Langlandscape.Comparison.xlambda\_finite}, \pathtt{Comparison/XLambda.lean}): for each $s\ge1$ it gives $\mathtt{Nonempty}$ of a $\Zp$-linear equivalence $e(\Zp\otimes G.X_s)\simeq e\,\Zp[S_s]$, with $e$ the ordinary part of $U_p\otimes1$ resp.\ of $U_p$. It is a corollary of \lean{Langlandscape.Comparison.xlambda\_finite\_iso} (under $\HypXchar$, a $\Zp$-linear isomorphism $\Psi:\Zp\otimes G.X_s\simeq\Zp[S_s]^0$ intertwining $U_p\otimes1$ and $U_p$). The Lean statement is only an existence (rank-level) statement: that the equivalence is \emph{induced by} $\varphi_s$ is not part of the \texttt{Nonempty} statement (it is for the definition \texttt{xlambdaOrd}; the open point R2-2 of the Lean audit). The inverse limit over $s$, parts (b)--(d) and the Galois and monodromy statements are not formalized (\leanplanned{}).
\end{quoted}

\begin{quoted}{Conjecture}{18.40}{159}{conj:s05-O2}{O2, a simplicial realisation; open}{chapters/s05/simplicial.tex}
There is a Hecke-equivariant map of finite simplicial (orbi-)sets, defined over $\Z[\frac1{6}]$ at most, from $\GDR_q(U_s)$ to a simplicial
model of the special fibre of $X_s$ at $q$ (or to a $q$-adic tubular neighbourhood), inducing $\varphi_s^{-1}$ on $H_1$. \stconj{};
equivalently (by Remark~18.38) a degree-shifting map between the two $S$-arithmetic complexes. The most promising concrete
form is the Deuring map on supersingular points, which exists only after reduction mod $q$; we expect no characteristic-zero Betti
version. \stopen.
\end{quoted}


\subsection{Stage 06: the integral Jacquet--Langlands defect}
Chapter: \mainpdf, pp.~134--162. On the case $q\equiv-1\bmod p$ with $\bar\rho$ unramified at $q$: Fabian lifted the
\verb|\provisional| marks on (F1) and (F1$'$) (15 November 2026); $N(\bar\rho)>1$ and non-squarefree $M$ ($\ell^2\mid M$ at a prime where $\bar\rho$ is unramified) are \stopen{} and are not claimed
(Prop.~23.2 and the note after the stage-06 boxes).
% SPDX-License-Identifier: CC-BY-4.0
% Copyright (C) 2026 Fabian Januszewski, Abhijit (the Langlandscape team)
% GENERATED by tools/extract_statements.py from mathematics/chapters/ -- do not edit by hand.

\begin{quoted}{Hypothesis}{20.1}{165}{hyp:s06-H}{standing hypotheses; seed \S9.1, corrected}{chapters/s06/s06-m2.tex}
\leavevmode
\begin{itemize}
\item[(H0)] $p\ge5$, $q\ne p$, $(M,pq)=1$, $N=qM$, $p\nmid\varphi(N)$. In particular $q\not\equiv1\bmod p$.
\item[(H0$^-$)] (\emph{new}; needed for the fixed-weight input Cor.~21.2 and for Prop.~24.1;
  Thm~21.5 needs (H0$^-$) \emph{or} (Lev); the inclusion of Prop.~22.1 uses neither)
  $q\not\equiv-1\bmod p$, or $\bar\rho_\m$ is ramified at $q$. Together with (H0): $p\nmid q^2-1$, or $\bar\rho_\m$ is ramified at $q$.
\item[(Lev)] (\emph{new}; alternative to (H0$^-$) in Thm~21.5) At the point $P$ of (P), of weight $k=k_P$, with
  $f=F_P$ (ordinary): either (F0) $M=1$; or (F1) $M>1$ is squarefree, $\bar\rho$ is unramified at every $\ell\mid M$ (so $N(\bar\rho)=1$),
  and $\ell\not\equiv-1\bmod p$ for every $\ell\mid M$; or (F1$'$) as (F1), but with some $\ell\mid M$, $\ell\equiv-1\bmod p$. These are the
  cases (F0), (F1), (F1$'$) of Prop.~23.2, in which its hypothesis (Lev$'_M$) (with (L3)) holds automatically. Status:
  Prop.~23.2 is \stpmc{} in all three cases, for (F1$'$) in our sign-fixed reading R4 of \cite[Thm~3.14]{KO} (ours, not
  \cite{KO}'s text). Not included: $N(\bar\rho)>1$ and non-squarefree $M_0$, which are \stopen{} (Remark~23.8); (L3) is the condition under
  which the method may extend.

\item[(H1)] $\T$ is Hida's full ordinary cuspidal Hecke algebra of tame level $\Gamma_0(N)$ over $\Lam$, in a fixed weight component
  $k_0\bmod(p-1)$. It is generated by all $T_\ell$ ($\ell\nmid Np$), $U_\ell$ ($\ell\mid Np$, in particular $U_q$ and $U_p$), and it
  is localised at a non-Eisenstein maximal ideal $\m$. $F$ is a primitive ordinary $q$-new branch, with $\lambda_F:\T\to I_F$, and
  $I_F$ is a normal local domain finite over $\Lam$. Then $I_F$ is $\Lam$-free: it is normal of dimension $2$, hence
  Cohen--Macaulay (Serre's criterion), hence free over the regular local ring $\Lam$ (Auslander--Buchsbaum). So the seed's
  ``hence $\Lam$-free'' is correct; it is used only through flatness of $R=\T\otimes_\Lam I_F$ over $I_F$.
\item[(H2)] $\bar\rho=\bar\rho_\m$ is absolutely irreducible on $G_{\Q(\sqrt{p^*})}$ and $p$-distinguished.
\item[(Im)] $p\mid\#\bar\rho(G_\Q)$. \emph{None of the cited theorems below uses (Im).} \cite[Thm~1.1]{KO} needs only (H2) (its
  condition (2) $=$ Assumption 2.1(TW)), and \cite[Rem.~2.10]{KO} says that the \cite{DFG} and \cite{Dim09} image conditions are interchangeable
  over $\Q$. We keep (Im) as in the seed, flagged as probably superfluous (Risk R4).
\item[(P)] Either (P1) $k_0\not\equiv2\bmod(p-1)$, and then $F$ has an arithmetic point $P$ of weight $k_P\in[3,p-1]$ with trivial
  character (a \emph{Kim--Ota point}). Or (P2) $k_0\equiv2$, $F$ has a weight-two point $P$ which is the ordinary
  $p$-stabilisation of a newform $f$ of level $N$, \emph{and} (\emph{new}) the weight-two ordinary Hecke algebra
  $\T^{\ord}_2(\Gamma_0(Np))_\m$ has no $p$-new eigenforms. The extra condition makes $\T/P\T\cong\T_2(\Gamma_0(N))_\m$,
  which is the algebra of \cite{KO}. Without it the control step (3) of the proof compares with an algebra that \cite{KO} does not treat.
\end{itemize}
\sthyp
\end{quoted}

\begin{quoted}{Theorem}{19.8}{163}{thm:s06-E}{the combinatorial half; seed Thm~E, hypotheses explicit}{chapters/s06/s06-m1.tex}
Let $S,\fa,\fb$ be as in Lemma~19.5. Put $S'=S/\fa$, and let $\lambda:S\to A$ factor through $\lambda':S'\to A$. Assume
\begin{itemize}
\item[(Gor$_S$)] $S$ is $A$-Gorenstein.
\end{itemize}
Let $\theta$ generate $S[\lambda]$ (Lemma~19.3), and let $x\in S'$ be its image. Then:
\begin{enumerate}[label=(\roman*)]
\item $\theta\in\fb$, $x\in S'[\lambda']$, $\eta(S,\lambda)=(\lambda'(x))$, and $\cc_q:=\lambda(\fb)=\{\psi(x):\psi\in S'^\vee\}$ (the
  \emph{order ideal} of $x$).
\item If $S'[\lambda']$ is free of rank one with a generator $\theta'$ that is unimodular in $S'$ (i.e.\ $S'[\lambda']$ is an
  $A$-direct summand of $S'$), then $\eta(S,\lambda)=\cc_q\cdot\eta(S',\lambda')$.
\item This holds in particular under (Gor$_{S'}$): $S'$ is finite free and $A$-Gorenstein. It also holds whenever $A$ is a DVR
  (\emph{no} hypothesis on $S'$). For $A$ normal, the identity therefore holds after localisation at every height-one prime.
\item In general $\cc=\mathrm{im}(\fb\to S')\cong\fb\cong S'^\vee$. So, if $S'$ is $A$-free (e.g.\ $S/\fb$ free, Lemma~19.5(c)), $\cc$ is a principal ideal of $S'$ if and only if
  (Gor$_{S'}$) holds.
\end{enumerate}
With $S=R$, $A=I_F$: under (Gor$_R$) and (Gor$_{R'}$), $\eta_F=\cc_q(F)\cdot\eta^B_F$, i.e.\ $(h_F)=\cc_q(F)$. \stproved{}
(Lean, namespace \pathtt{Langlandscape.CommAlg.ThmE}: (i) \lean{Langlandscape.CommAlg.ThmE.theta\_mem\_b}, \pathtt{mk\_theta\_mem\_ann}, \pathtt{eta\_eq\_span\_lambda'}, \lean{Langlandscape.CommAlg.ThmE.c\_eq\_orderIdeal};
(ii) \lean{Langlandscape.CommAlg.ThmE.eta\_eq\_c\_mul\_of\_unimodular};
(iii) Gorenstein half \lean{Langlandscape.CommAlg.ThmE.unimodular\_of\_isGorenstein}, \lean{Langlandscape.CommAlg.ThmE.eta\_eq\_c\_mul\_of\_isGorenstein}, DVR half \lean{Langlandscape.CommAlg.ThmE.eta\_eq\_c\_mul\_of\_dvr} (\pathtt{CommAlg/DVR.lean}, with \pathtt{exists\_unimodular\_of\_dvr});
(iv) \lean{Langlandscape.CommAlg.ThmE.c\_principal\_iff\_isGorenstein}.)
\stformalized{} for (i), (ii), (iii) (both halves) and (iv) for $S'$ $A$-free, under the hypotheses of Lemma~19.5 as the structure \texttt{NewOldData} (same qualifiers as there). \emph{Not formalized:} the localisation clause of (iii) for $A$ normal (height-one primes), the isomorphism $\mathrm{im}(\fb\to S')\cong\fb\cong S'^\vee$ as a standalone statement for general $S'$, and the final specialisation $S=R$, $A=I_F$ (the abstract statement is formalized, its instantiation to the Hecke rings is not).
\end{quoted}

\begin{quoted}{Proposition}{22.1}{174}{prop:s06-galois}{Galois comparison; seed Prop.~galois}{chapters/s06/s06-m4.tex}
Assume (H0), (H1), (H2). Then $\tq_q(F)\subseteq\cc_q(F)$. \stpmc{}
(Lean: Given \pathtt{Hyp.ResidualBurnside}, \pathtt{Hyp.FreeRankTwo} is equivalent to the inclusion $\tq_q\subseteq\cc_q$; both directions are proved in Lean (\pathtt{Defect.galois}, \pathtt{Defect.freeRankTwo\_of\_le}, packaged as \pathtt{Defect.freeRankTwo\_iff\_le}). The Galois half is therefore not formalized independently: its arithmetic content (that the old-form representation $\rho_M$ exists at level $M$, steps 1--3 of Prop.~22.1) is the cited hypothesis. \pathtt{Hyp.ResidualBurnside} is discharged only given absolute irreducibility, (H2).)
\par\smallskip\noindent{\scriptsize\itshape [Status paragraph following the statement in the chapter:]}\par
\noindent\emph{Cited inputs.}
(1) $H$ is free of rank two over $\T$, and $H(M)$ is free of rank two over $\T(M)_\m$ (tame level $M$):
\cite[\S7.2, p.~68]{Hs} \lit{hsieh\_2017\_hida-families-and-p-adic-triple-product-l-functions\_arxiv}{\S7.2, proof of Lemma 7.4, p.~68}
\stcited, citing \cite[Cor.~1, 2, p.~482]{Wi95}
\lit{wiles\_1995\_modular-elliptic-curves-and-fermats-last-theorem}{Thm 2.1, Cor.~1, Cor.~2, p.~482}, under (CR) (which (H2)
implies, Remark~20.2). The Galois action on $H$ (resp.\ $H(M)$) is unramified outside $Np$ (resp.\ $Mp$), and
$\tr(\Frob_\ell\mid H)=T_\ell$ in $\T$ for $\ell\nmid Np$ (Eichler--Shimura; for each classical specialisation this is
\lit{darmon-diamond-taylor\_1997\_fermats-last-theorem\_preprint}{Thm 3.1(a), p.~85}) \stcited.
(2) $H(M)$ is unramified at $q$, since $q\nmid Mp$ ($X_1(Mp^r)$ has good reduction at $q$; for each classical specialisation this is
again \lit{darmon-diamond-taylor\_1997\_fermats-last-theorem\_preprint}{Thm 3.1(a), p.~85}) \stknown.
(3) Chebotarev: for a Galois extension of $\Q$ unramified outside a finite set $S$, the Frobenius elements at primes $\notin S$ are
dense \lit{darmon-diamond-taylor\_1997\_fermats-last-theorem\_preprint}{Thm 2.3, p.~52} \stknown.
\end{quoted}

\begin{quoted}{Theorem}{21.5}{172}{thm:s06-defect}{$\Lam$-adic defect formula; seed Thm~defect, corrected}{chapters/s06/s06-m3.tex}
Assume (H0), (H1), (H2), (P) (with (P2) in its corrected form), and (H0$^-$) or (Lev) (Hyp.~20.1). Then
\[(h_F)=\cc_q(F)=\tq_q(F)\qquad\text{as principal ideals of }I_F.\]
\stpmc{} under (H0$^-$). Under (Lev) (this is where $q\equiv-1\bmod p$ with $\bar\rho$ unramified at $q$ is covered): \stpmc{} for $M=1$
and for $M>1$ in the cases (F1), (F1$'$) of Prop.~23.2 (for (F1$'$) in our sign-fixed reading of \cite[Thm~3.14]{KO}, not
\cite{KO}'s text). For $q\equiv-1\bmod p$ with $\bar\rho$ unramified at $q$ and outside (Lev) ($N(\bar\rho)>1$, or $\ell^2\mid M$ at a prime $\ell$ where $\bar\rho$ is unramified) the theorem is \stopen{} (Remark~25.1(i), Remark~23.8).

 Cited inputs: Hida freeness
and control \cite[Thm~3.1, Cor.~3.2]{Hi86b} \lit{hida\_1986\_iwasawa-modules-attached-to-congruences-of-cusp-forms}{Thm 3.1, Cor.~3.2
(3.3b), p.~250}; (Gor$_R$) and (Gor$_{R'}$), with freeness of $H$ (\cite{Hs} as above, resting on \cite{Wi95}, \cite{MW86}, \cite[Cor.~8.11]{Hel}
\lit{helm\_2004\_maps-between-modular-jacobians-and-jacobians-of-shimura-curves\_arxiv}{Cor.~8.11, p.~31}); the $\Lam$-adic JL
(\pathtt{if:s01-targets}); and Thm~21.1 at the point $P$ of (P) under (H0$^-$), or Prop.~23.2 at $P$ under
(Lev) (resting on \cite[Thm~2.9, Lemma~2.4]{KO} for $M=1$, \cite[Thm~3.14]{KO} for (F1), (F1$'$)). (Im) is not used.
\end{quoted}

\begin{quoted}{Corollary}{21.6}{173}{cor:s06-allweights}{all weights of the branch; seed Cor.~allweights}{chapters/s06/s06-m3.tex}
Under the hypotheses of Thm~21.5, let $P'$ be any arithmetic point of weight $k'\ge3$ with trivial character
(including $k'\ge p$), with $q$-new $f'=F_{P'}$ ($F_{P'}$ is the $p$-stabilisation of a newform of level $N$, $F$ being primitive of tame level $N$). Then $\cc_q(f')=(\varpi_{P'}^{t_{f'}(q)})$ (in $\widetilde\cO_{P'}$) and
$\ord_{\varpi_{P'}}\eta_{f'}(N)-\ord_{\varpi_{P'}}\eta_{f'}(M,q)=t_{f'}(q)$, for the ring-theoretic congruence ideals of $\T_{k'}(N)_\m$ and of its
$q$-new quotient. \stpmc{} (with the status of the case of Thm~21.5 that is used; under (Lev), case (F1$'$), in our reading of \cite[Thm~3.14]{KO}).
\end{quoted}

\begin{quoted}{Proposition}{23.2}{176}{prop:s06-levM}{the inclusion without (H0$^-$); the seed's step (4)}{chapters/s06/s06-m5.tex}
Let $f$ be as above, with (H0) and (Lev$'_M$) (Hyp.~23.1), and assume that $f$ is $p$-ordinary, i.e.\ $a_p(f)\in\cO^\times$
(automatic for $f=F_P$ in Thm~21.5). Then $\cc_q(f)\subseteq(\varpi^{t_f(q)})$ in each of the following cases, in which
(Lev$'_M$) holds automatically, $N(\bar\rho)=1$, $M_0=M$ and $\Sigma^+=P_{\bar\rho}=\emptyset$:
\begin{itemize}
\item[(F0)] $M=1$. \stpmc{} (Remark~23.10).
\item[(F1)] $M>1$ squarefree, $\bar\rho$ unramified at every $\ell\mid M$, and $\ell\not\equiv-1\bmod p$ for every $\ell\mid M$.
  \stpmc
\item[(F1$'$)] as (F1), but some $\ell\mid M$ has $\ell\equiv-1\bmod p$. \stpmc{} in the sign-fixed reading of \cite[Thm~3.14]{KO}
  (Remark~23.4), which is our reading, not \cite{KO}'s text (F2 part B and audit part 2: confirmed as a reading).

\end{itemize}
No condition on $q\bmod p$ beyond (H0) is used; in particular $q\equiv-1\bmod p$ is allowed. For $N(\bar\rho)>1$ (under (Lev$'_M$), with
or without (L3)) the inclusion is \stopen: the proof below does not apply when $\bar\rho$ is not twist-minimal (counterexample B2 of F2,
Remark~23.9), and the inclusion itself is neither proved nor disproved there.
\end{quoted}

\begin{quoted}{Proposition}{24.1}{184}{prop:s06-betti}{Betti incarnation of $\cc_q$}{chapters/s06/s06-m7.tex}
Let $l=p>k\ge4$ (the range of \pathtt{if:s03-amalgam}), $q\not\equiv-1\bmod p$, and let $\m$ be non-Eisenstein for $\T_k(M)$ acting on
$H^1(\Gamma_0(M),V_k)$, with $H^1(\Gamma_0(M),V_k)_\m$ free of rank two over $\T_k(M)_\m$. Then:
\begin{enumerate}[label=(\alph*)]
\item $\coker(\Psi^*\Psi)_\m\cong\big(\T_k(M)_\m/(\delta)\big)^{\oplus2}$, where $\delta:=\det(\Psi^*\Psi)=(q+1)^2-q^{2-k}T_q^2$.
  Since $p\nmid2(q+1)$, $(\delta)=(T_q-\epsilon(q+1)q^{(k-2)/2})$ for the unique sign $\epsilon$ with this element in $\m$.
\item $\T_k(N)_{\m}/(\fa+\fb)\cong\T_k(M)_\m/(T_q-\epsilon(q+1)q^{(k-2)/2})$ (Ribet's level-raising module, with Ihara's lemma).
\item If $\rk_\cO\T'_{k,\m}=1$, then $\ord_\varpi\cc_q(f)=\tfrac12\,\mathrm{length}_\cO\coker(\Psi^*\Psi)_\m$. In general only
  $\ord_\varpi\cc_q(f)\le\tfrac12\mathrm{length}$ holds.
\end{enumerate}
(a) \stproved{} (linear algebra, Lemma~24.2). (b) \stpmc{} (Ihara's lemma in weight $k$ and the cohomological
congruence module: \cite[Lemma~3.2 and \S2]{Dia91}, \cite[\S3]{Dia89}). (c) \stpmc{} (from (a), (b)).
Here: in (a), if neither sign gives an element of $\m$, then $\delta\in\T_k(M)_\m^\times$ and both sides are $0$ for either sign;
in (b), $\T_k(N)_\m$ is the localisation of the weight-$k$ Hecke algebra of $\Gamma_0(N)$ (with $U_q$; $N=qM$) at the maximal ideal $\m_N$
generated by $\m$ (on the $T_\ell$, $\ell\nmid Np$, and $U_\ell$, $\ell\mid M$) and $U_q-\epsilon c$, and $\fa,\fb$ are its $q$-new and
$q$-old kernels; in (c), $f$ is a $q$-new eigenform in $\m_N$, $\T'_{k,\m}$ is the $q$-new quotient of $\T_k(N)_\m$, $\cc_q(f)$ is as in
Def.~20.5, and (H0) is assumed. \provisional{S06-M7: the degenerate case ``no sign'' in (a) and the maximal ideal
$\m_N$ of (b) are made explicit, and (H0) ($p\nmid\varphi(N)$, $q\not\equiv1$) is used in (b); statement otherwise unchanged.}
\end{quoted}

\begin{quoted}{Conjecture}{25.2}{187}{conj:s06-several}{several ramified primes; seed \S10, last paragraph}{chapters/s06/s06-m8.tex}
Let $D=q_1\cdots q_r$ ($r$ odd), with each $q_i$ satisfying (H0), (H0$^-$) and $F$ $D$-new. Let
$\T=\T^{(0)}\twoheadrightarrow\T^{(1)}\twoheadrightarrow\cdots\twoheadrightarrow\T^{(r)}=\T^{D\text{-new}}\cong\T^{\ord}_{B_D}$ be the
$q_1\cdots q_i$-new quotients, and let $\cc^{(i)}_{q_i}(F)$ be the old/new congruence ideal of $\T^{(i-1)}\twoheadrightarrow\T^{(i)}$ at
$q_i$. Then $(h_F)=\prod_i\cc^{(i)}_{q_i}(F)=\prod_i\tq_{q_i}(F)$. \stconj{} If every $\T^{(i)}_\m$ is Gorenstein, the first
equality follows by iterating Thm~19.8 \stsketch. Whether the intermediate quotients are Gorenstein is \stopen{} (Hsieh
covers $i=0$ and $i=r$). The second equality would follow from Prop.~22.1 at each step and \cite[Thm~1.1]{KO} with
$N^-=D$, by the squeeze of step (5) \stsketch. Note that $q_i$ is $q_1\cdots q_{i-1}$-new-ramified in the intermediate algebras,
which needs the relative version of Lemma~20.4 at each step.
\end{quoted}

\noindent\emph{Note.} The status remark for Prop.~23.2 (Remark~23.8 of \mainpdf, p.~156) records: (F0) \stpmc; (F1) \stpmc; (F1$'$) \stpmc{} in our sign-fixed reading R4 of \cite[Thm~3.14]{KO},
which is our reading, not \cite{KO}'s text; $N(\bar\rho)>1$ \stopen. There is no numerical test of (F1) or (F1$'$).

\subsection{All results of the chapters}
{\small
% SPDX-License-Identifier: CC-BY-4.0
% Copyright (C) 2026 Fabian Januszewski, Abhijit (the Langlandscape team)
% GENERATED by tools/extract_statements.py from mathematics/chapters/ -- do not edit by hand.

\begin{longtable}{@{}>{\raggedright}p{0.13\textwidth} >{\raggedright}p{0.25\textwidth} >{\raggedright}p{0.30\textwidth} >{\raggedright\arraybackslash}p{0.25\textwidth}@{}}
\caption{Every theorem, proposition, corollary and conjecture of \texttt{mathematics/main.pdf} (number, page), with the status markers that occur inside its statement and, marked ``after the statement'', those of a status paragraph printed directly after it in the chapter (generated; \texttt{\textbackslash lean} refs = number of Lean cross-references in the statement). Mixed markers refer to different clauses; the statement in main.pdf is authoritative. A ``--'' means that no marker was found in either place, not that the result is proved. Rows are in the numerical order of main.pdf within each stage.}\label{tab:allresults}\\
\toprule no.\ (p.) & label & title & status markers\\\midrule\endfirsthead
\toprule no.\ (p.) & label & title & status markers\\\midrule\endhead
\multicolumn{4}{@{}l}{\textbf{Stage 01}}\\
Prop. 1.4 (p.~28) & \texttt{\footnotesize prop:s01-eichler} & Eichler orders & \stknown{}\\
Prop. 1.6 (p.~29) & \texttt{\footnotesize prop:s01-finite} & finiteness & \stknown{}, \stformal{}; 2 \texttt{\textbackslash lean} refs\\
Theo. 2.1 (p.~31) & \texttt{\footnotesize thm:s01-mass} & Eichler mass formula; \texttt{if:s01-mass} & \stcited{}, \stformal{}, \stproved{}, \stpmc{}; 1 \texttt{\textbackslash lean} refs\\
Prop. 3.3 (p.~33) & \texttt{\footnotesize prop:s01-adjoint} & adjoints; equivariance of $\nu$ & \stcited{}, \stproved{}, \stformal{}; 8 \texttt{\textbackslash lean} refs\\
Prop. 4.4 (p.~39) & \texttt{\footnotesize prop:s01-relations} & relations & \stproved{}, \stformal{}, \stknown{}; 4 \texttt{\textbackslash lean} refs\\
Theo. 6.1 (p.~48) & \texttt{\footnotesize thm:s01-JLclassical} & classical JL, Brandt form; \pathtt{if:s01-targets}; seed Thm.~3.1 & \stcited{}; after the statement: \stformal{}\\
Theo. 6.2 (p.~49) & \texttt{\footnotesize thm:s01-chenevier} & $p$-adic JL; Chenevier & \stcited{}\\
Theo. 6.3 (p.~49) & \texttt{\footnotesize thm:s01-lambdaJL} & $\Lambda$-adic Hecke JL & \stcited{}, \textsf{\small[provisional note]}; 1 \texttt{\textbackslash lean} refs; after the statement: \stformal{}\\
\multicolumn{4}{@{}l}{\textbf{Stage 02}}\\
Theo. 7.7 (p.~54) & \texttt{\footnotesize thm:s02-SA} & strong approximation for $B^1$ at $S=\{\infty,q\}$ & after the statement: \stcited{}, \stproved{}\\
Prop. 7.9 (p.~55) & \texttt{\footnotesize prop:s02-BG} & SA for one level as a finite check & after the statement: \stproved{}, \stcited{}, \stformal{}\\
Prop. 7.10 (p.~57) & \texttt{\footnotesize prop:s02-quot} & quotient graph; seed Prop.~quot & \stcited{}; after the statement: \stproved{}, \stpmc{}\\
Prop. 7.14 (p.~61) & \texttt{\footnotesize prop:s02-Mhomprops} & properties of Def.~7.13 & after the statement: \stproved{}, \stformal{}\\
Prop. 7.16 (p.~64) & \texttt{\footnotesize prop:s02-duality} & the two pairings & \stformal{}; 2 \texttt{\textbackslash lean} refs; after the statement: \stproved{}\\
Prop. 7.18 (p.~65) & \texttt{\footnotesize prop:s02-coh} & seed Prop.~coh & after the statement: \stcited{}, \stpmc{}, \stproved{}, \stformal{}\\
Prop. 7.19 (p.~66) & \texttt{\footnotesize prop:s02-harm} & harmonic cochains; seed Prop.~harm & after the statement: \stproved{}, \stpmc{}, \stformal{}\\
Prop. 7.20 (p.~67) & \texttt{\footnotesize prop:s02-free} & the free case & after the statement: \stpmc{}, \stproved{}, \stformal{}\\
Prop. 7.22 (p.~68) & \texttt{\footnotesize prop:s02-indefprops} & properties of the tree model & after the statement: \stpmc{}, \stproved{}, \stformal{}\\
Prop. 7.31 (p.~73) & \texttt{\footnotesize prop:s02-sa-step-d} & tex L02.SA(d) & \stpmc{}, \stcited{}\\
Theo. 7.34 (p.~74) & \texttt{\footnotesize thm:s02-sa-main} & Thm.~7.7, proof & \stpmc{}\\
\multicolumn{4}{@{}l}{\textbf{Stage 03}}\\
Prop. 8.3 (p.~76) & \texttt{\footnotesize prop:s03-phi} & two-vertex graphs; seed Prop.~phi & \stproved{}, \stformal{}; 4 \texttt{\textbackslash lean} refs\\
Prop. 8.8 (p.~79) & \texttt{\footnotesize prop:s03-mazur} & Mazur check; seed Example mazur & \stpmc{}, \stformal{}; 8 \texttt{\textbackslash lean} refs\\
Prop. 9.5 (p.~83) & \texttt{\footnotesize prop:s03-phieis} & $\Phicomb$ is Eisenstein & \stproved{}, \stformal{}, \stcited{}; 4 \texttt{\textbackslash lean} refs\\
Prop. 9.7 (p.~85) & \texttt{\footnotesize prop:s03-LRdef} & definite level-raising operator & \stproved{}, \stformal{}, \stcited{}; 11 \texttt{\textbackslash lean} refs\\
Theo. 10.2 (p.~89) & \texttt{\footnotesize thm:s03-ribet} & definite Ribet sequence as a vertex collapse; seed Thm ribet & \stproved{}, \stpmc{}, \stcited{}, \stformal{}, \leanplanned{}; 14 \texttt{\textbackslash lean} refs\\
Coro. 10.5 (p.~90) & \texttt{\footnotesize cor:s03-ihara} & definite Ihara lemma, (i)--(ii); seed Cor.~ihara & \stproved{}, \stpmc{}, \stformal{}; 4 \texttt{\textbackslash lean} refs\\
Theo. 11.3 (p.~95) & \texttt{\footnotesize thm:s03-iharaH0} & Ihara (iii): $H_0$ and its Hecke action & \stpmc{}, \sttoprove{}, \stsketch{}, \stproved{}, \stformal{}, \leanplanned{}; 5 \texttt{\textbackslash lean} refs\\
Coro. 11.5 (p.~98) & \texttt{\footnotesize cor:s03-eis} & Eisenstein support & \stpmc{}, \stcited{}, \leanplanned{}\\
Prop. 12.2 (p.~101) & \texttt{\footnotesize prop:s03-amalgam} & amalgam & \stpmc{}, \stcited{}\\
Prop. 12.3 (p.~101) & \texttt{\footnotesize prop:s03-MV} & Mayer--Vietoris; seed \S6.4 & \stpmc{}, \stcited{}\\
Prop. 12.6 (p.~104) & \texttt{\footnotesize prop:s03-PsiPsi} & $\Psi^*\Psi$ & \stpmc{}, \sttoprove{}, \stcomp{}\\
Conj. 12.10 (p.~107) & \texttt{\footnotesize conj:s03-iharaGL2} & Ihara lemma for $\GL_2$, Betti form & \stexp{}, \stcited{}\\
\multicolumn{4}{@{}l}{\textbf{Stage 04}}\\
Prop. 14.2 (p.~112) & \texttt{\footnotesize prop:s04-freeA} & freeness at finite level; seed Prop.~free (a) & \stproved{}\\
Theo. 14.4 (p.~114) & \texttt{\footnotesize thm:s04-control} & $\Lambda$-freeness and control; seed Prop.~free (b); \texttt{if:s04-Mord} & \stexp{}, \stsketch{}, \stproved{}, \stpmc{}, \stcited{}, \stformal{}, \textsf{\small[provisional note]}; 7 \texttt{\textbackslash lean} refs\\
Theo. 15.1 (p.~116) & \texttt{\footnotesize thm:s04-lambdaribet} & $\Lambda$-adic definite Ribet sequence; seed Thm.~lambdaribet; \pathtt{if:s04-lambdaribet} & \stproved{}, \stformal{}, \textsf{\small[provisional note]}; 5 \texttt{\textbackslash lean} refs\\
Theo. 16.6 (p.~122) & \texttt{\footnotesize thm:s04-Gor} & Gorenstein and freeness; \texttt{if:s04-Mord}, second half & \stexp{}; after the statement: \stcited{}, \stproved{}\\
Theo. 16.7 (p.~123) & \texttt{\footnotesize thm:s04-etaB} & $\eta^B_F$ as a self-pairing; \texttt{if:s04-etaB} & after the statement: \stcited{}\\
Prop. 17.1 (p.~124) & \texttt{\footnotesize prop:s04-etaprinciple} & $\eta$-principle; \pathtt{thm:s04-etaB}, \texttt{lem:s06-S} & \stproved{}, \stformal{}; 4 \texttt{\textbackslash lean} refs\\
Prop. 18.1 (p.~125) & \texttt{\footnotesize prop:s04-ekillsVk} & $e$ kills $H_0$ with $\Vk$-coefficients & after the statement: \sttoprove{}\\
\multicolumn{4}{@{}l}{\textbf{Stage 05}}\\
Prop. 18.6 (p.~128) & \texttt{\footnotesize prop:s05-sign} & $\varepsilon=+1$; exponent $+1$ & \stpmc{}, \stproved{}, \stknown{}, \sttoprove{}, \sthyp{}, \textsf{\small[provisional note]}\\
Prop. 18.13 (p.~137) & \texttt{\footnotesize prop:s05-signe} & Prop.~sign (e): higher weight & \stpmc{}, \stexp{}\\
Prop. 18.18 (p.~141) & \texttt{\footnotesize prop:s05-AL} & Atkin--Lehner variant & \stpmc{}, \stcited{}, \stexp{}, \stsketch{}, \textsf{\small[provisional note]}\\
Prop. 18.20 (p.~144) & \texttt{\footnotesize prop:s05-Xlambda} & $\Lam$-adic version; seed Prop.~Xlambda & \stproved{}, \stpmc{}, \stopen{}, \stsketch{}, \stconj{}, \textsf{\small[provisional note]}; after the statement: \stformal{}, \leanplanned{}\\
Prop. 18.24 (p.~150) & \texttt{\footnotesize prop:s05-Xiso} & $\HypXiso$ holds non-canonically & \stpmc{}, \stformal{}, \stsketch{}; 1 \texttt{\textbackslash lean} refs\\
Prop. 18.25 (p.~151) & \texttt{\footnotesize prop:s05-subquot} & $q$-new sub- versus quotient lattice; seed Prop.~subquot; \pathtt{if:s05-subquot} & \stproved{}, \stformal{}; 2 \texttt{\textbackslash lean} refs\\
Prop. 18.28 (p.~152) & \texttt{\footnotesize prop:s05-nochain} & no tree-only chain; corrected form of seed Prop.~noChain & \stproved{}, \stpmc{}\\
Theo. 18.30 (p.~153) & \texttt{\footnotesize thm:s05-G} & (G1)--(G4) at $\Gamma_0$-type level & \stcited{}; after the statement: \stcited{}\\
Conj. 18.40 (p.~159) & \texttt{\footnotesize conj:s05-O2} & O2, a simplicial realisation; open & \stconj{}, \stopen{}\\
\multicolumn{4}{@{}l}{\textbf{Stage 06}}\\
Theo. 19.8 (p.~163) & \texttt{\footnotesize thm:s06-E} & the combinatorial half; seed Thm~E, hypotheses explicit & \stproved{}, \stformalized{}; 7 \texttt{\textbackslash lean} refs\\
Theo. 21.1 (p.~170) & \texttt{\footnotesize thm:s06-KO} & Kim--Ota & \stcited{}\\
Coro. 21.2 (p.~170) & \texttt{\footnotesize cor:s06-fixed} & fixed-weight level lowering mod $\varpi^t$; replaces seed step (4) & \stpmc{}\\
Theo. 21.5 (p.~172) & \texttt{\footnotesize thm:s06-defect} & $\Lam$-adic defect formula; seed Thm~defect, corrected & \stpmc{}, \stopen{}\\
Coro. 21.6 (p.~173) & \texttt{\footnotesize cor:s06-allweights} & all weights of the branch; seed Cor.~allweights & \stpmc{}\\
Prop. 22.1 (p.~174) & \texttt{\footnotesize prop:s06-galois} & Galois comparison; seed Prop.~galois & \stpmc{}; after the statement: \stcited{}, \stknown{}\\
Prop. 23.2 (p.~176) & \texttt{\footnotesize prop:s06-levM} & the inclusion without (H0$^-$); the seed's step (4) & \stpmc{}, \stopen{}\\
Prop. 24.1 (p.~184) & \texttt{\footnotesize prop:s06-betti} & Betti incarnation of $\cc_q$ & \stproved{}, \stpmc{}, \textsf{\small[provisional note]}\\
Conj. 25.2 (p.~187) & \texttt{\footnotesize conj:s06-several} & several ramified primes; seed \S10, last paragraph & \stconj{}, \stsketch{}, \stopen{}\\
\bottomrule\end{longtable}
}
